% ============================================================
%  Formelsammlung Elektrodynamik  —  Hauptdatei
%  TH Nürnberg, Studiengang MSY
% ============================================================
\documentclass[a4paper,10pt]{article}
\usepackage[utf8]{inputenc}
\usepackage[T1]{fontenc}
\usepackage[ngerman]{babel}
\usepackage{amsmath,amssymb,mathtools}
\usepackage{xcolor}
\usepackage{geometry}
\usepackage{array}
\usepackage{multirow}
\usepackage{microtype}
\usepackage{needspace}
\geometry{top=10mm,bottom=14mm,left=10mm,right=10mm,includefoot}
\setlength{\parskip}{0pt}
\setlength{\parindent}{0pt}
\renewcommand{\baselinestretch}{0.9}
\setlength{\emergencystretch}{2em}
\hfuzz=3pt
\setlength{\fboxsep}{0.18em}
\setlength{\fboxrule}{0.3pt}

% ---- Farbschema: eine Farbe pro EDy-Skript ---------------
\colorlet{cKap0}{gray!80!black}    % EDy0: Übersicht / Einteilung
\colorlet{cKap1}{blue}             % EDy1: Einteilung EM-Felder
\colorlet{cKap2}{teal}             % EDy2: Elektrostatisches Feld
\colorlet{cKap3}{violet}           % EDy3: Magnetisches Feld
\colorlet{cKap4}{orange!90!red}    % EDy4: Stationäres Strömungsfeld
\colorlet{cKap5}{olive}            % EDy5: Zeitveränderliches Strömungsfeld
\colorlet{cKap6}{purple}           % EDy6: Elektromagnetische Wellen
\colorlet{cKap7}{cyan!70!black}    % EDy7: Wellenleiter
\colorlet{cKap8}{red!70!black}     % EDy8: Antennen

% ---- Bereichskopf ----------------------------------------
\newcommand{\bereich}[2][cKap1]{%
  \needspace{4\baselineskip}%
  \vspace{0.5em}%
  \noindent\colorbox{#1!10}{%
    \makebox[\dimexpr\linewidth-2\fboxsep\relax][l]{%
      \small\bfseries\color{#1!80!black}\strut\enspace #2}}%
  \nopagebreak%
  \vspace{0.1em}\par\noindent}

% ---- Trennlinie ------------------------------------------
\newcommand{\bereichende}{%
  \vspace{0.2em}%
  {\color{gray!50}\hrule height 0.3pt}%
  \vspace{0.2em}}

% ---- Drei-Spalten-Layout ---------------------------------
\newenvironment{formelreihe}{%
  \noindent\footnotesize%
  \setlength{\tabcolsep}{0pt}%
  \renewcommand{\arraystretch}{0.85}%
  \begin{tabular}{@{}p{0.32\linewidth}@{\hspace{0.02\linewidth}}%
                     p{0.32\linewidth}@{\hspace{0.02\linewidth}}%
                     p{0.32\linewidth}@{}}%
}{\end{tabular}\par}

% ---- Formel-Block ----------------------------------------
\newcommand{\formelblock}[2]{%
  \parbox[t]{\linewidth}{%
    {\scriptsize\scshape\color{gray!65!black} #1}\\[0.03em]%
    $\displaystyle #2$}}

% ---- Beispiel-Block (nur Spalte 3) -----------------------
\newcommand{\beispielblock}[2]{%
  \fcolorbox{orange!50}{yellow!15}{%
    \parbox[t]{\dimexpr\linewidth-2\fboxsep-2\fboxrule\relax}{%
      {\scriptsize\scshape\color{orange!80!black} Bsp: #1}\\[0.03em]%
      $\displaystyle #2$}}}

% ---- Info-Block ------------------------------------------
\newcommand{\infoblock}[1]{%
  \parbox[t]{\linewidth}{\scriptsize\itshape #1}}

% ==========================================================
\begin{document}
\begin{center}
  {\large\bfseries Formelsammlung Elektrodynamik — TH Nürnberg MSY}
\end{center}
\vspace{0.3em}

% ---- EDy0/EDy1: Einteilung der elektromagnetischen Felder ----
\bereich[cKap0]{Einteilung der elektromagnetischen Felder}
\begin{formelreihe}
  \parbox[t]{\linewidth}{%
    \formelblock{Zeitlich konstant $\partial/\partial t=0$}{\begin{aligned}
      &\text{Elektrostatik: } \vec{H},\vec{B}=0\\[0.1em]%
      &\text{Magnetostatik: } \vec{E},\vec{D}=0\\[0.1em]%
      &\text{Stationaeres Stroemungsfeld}
    \end{aligned}}\\[0.25em]%
    \formelblock{Langsam zeitveraenderl.\ $\partial/\partial t\neq0$}{\begin{aligned}
      &\tfrac{\partial\vec{B}}{\partial t}\approx 0\text{: quasistat.\ el.\ Feld}\\[0.1em]%
      &\tfrac{\partial\vec{D}}{\partial t}\approx 0\text{: quasistat.\ mag.\ Feld}
    \end{aligned}}} &
  \parbox[t]{\linewidth}{%
    \formelblock{Schnell zeitveraenderlich}{\text{Elektromagnetisches Feld}}\\[0.25em]%
    \infoblock{%
      \textbf{Materialgleichungen:}\\
      $\vec{D}=\varepsilon\vec{E}$,\;
      $\varepsilon=\varepsilon_0\varepsilon_r$\\[0.15em]
      $\vec{B}=\mu\vec{H}$,\;
      $\mu=\mu_0\mu_r$\\[0.15em]
      $\varepsilon_0=8{,}854\cdot10^{-12}\,\mathrm{As/(Vm)}$\\
      $\mu_0=4\pi\cdot10^{-7}\,\mathrm{Vs/(Am)}$}} &
  \beispielblock{Lichtgeschwindigkeit}{c=\tfrac{1}{\sqrt{\varepsilon_0\mu_0}}
    \approx 3\cdot10^8\,\tfrac{\mathrm{m}}{\mathrm{s}}} \\
\end{formelreihe}
\bereichende

% ---- EDy2: Elektrostatisches Feld ----------------------------

% 2.1 / 2.2  Coulomb & Elektrische Feldstärke
\bereich[cKap2]{2.1 Coulombsches Gesetz \& 2.2 Elektrische Feldstärke $\vec{E}$}
\begin{formelreihe}
  \parbox[t]{\linewidth}{%
    \formelblock{Coulomb-Kraft (Punktladungen)}{%
      \vec{F}=\tfrac{1}{4\pi\varepsilon}\cdot\tfrac{Q_1\cdot Q_2}{r^2}\cdot\vec{e}_r}\\[0.25em]%
    \formelblock{Elektrische Feldstärke}{%
      \vec{E}=\tfrac{\vec{F}}{Q_p}=\tfrac{1}{4\pi\varepsilon}\cdot\tfrac{Q}{r^2}\cdot\vec{e}_r}\\[0.25em]%
    \infoblock{Richtung $\vec{E}$ = Richtung von $\vec{F}$ auf positive Probeladung $Q_p$.\\
      $[\vec{E}]=1\,\tfrac{\text{N}}{\text{As}}=1\,\tfrac{\text{V}}{\text{m}}$}} &
  \parbox[t]{\linewidth}{%
    \formelblock{Verschiebungsdichte (materialunabh.)}{%
      \vec{D}=\varepsilon\cdot\vec{E}\;,\quad
      \left[\vec{D}\right]=\tfrac{\text{As}}{\text{m}^2}}\\[0.25em]%
    \formelblock{Punktladung $Q$}{%
      \vec{D}=\tfrac{Q}{4\pi r^2}\cdot\vec{e}_r}\\[0.25em]%
    \formelblock{Hüllfläche $A$}{%
      Q=\oint_A\!\vec{D}\,\mathrm{d}\vec{A}}\\[0.25em]%
    \formelblock{Divergenz (Gauss)}{%
      \operatorname{div}\vec{D}=\tfrac{\mathrm{d}Q}{\mathrm{d}V}=\rho}} &
  \parbox[t]{\linewidth}{%
    \beispielblock{Feldlinienbild}{%
      \begin{aligned}
        &\text{Richtung}\ \vec{E}=\text{Tangente}\\
        &\text{Betrag}\ |\vec{E}|=\text{Dichte}\\[0.2em]
        &\oint_A\!\vec{E}\,\mathrm{d}\vec{A}\neq0\;\Rightarrow\text{nicht quellfrei}
      \end{aligned}}\\[0.25em]
    \grafik[scale=0.6]{%
      \foreach \k in {0.32,0.72,1.2}{%
        \draw[fld] (-0.7,0) .. controls (-0.28,\k) and (0.28,\k) .. (0.7,0);
        \draw[fld] (-0.7,0) .. controls (-0.28,-\k) and (0.28,-\k) .. (0.7,0);}
      \draw[fld] (-0.55,0)--(0.6,0);
      \node[plus] at (-0.7,0){$+$};
      \node[minus] at (0.7,0){$-$};}} \\[0.3em]
\end{formelreihe}
\bereichende

% 2.3  Spannung und Potential
\bereich[cKap2]{2.3 Spannung und Potential}
\begin{formelreihe}
  \parbox[t]{\linewidth}{%
    \formelblock{Arbeit (homogenes Feld)}{%
      W=Q\cdot E\cdot s\cdot\cos\alpha}\\[0.25em]%
    \formelblock{Homogenfeld: Plattenkond.}{%
      E=\tfrac{U_{12}}{s_{12}}=\tfrac{U}{d}}\\[0.25em]%
    \formelblock{Allgemein (Linienintegral)}{%
      U_{12}=\int_1^2\!\vec{E}\,\mathrm{d}\vec{s}=-U_{21}}} &
  \parbox[t]{\linewidth}{%
    \formelblock{Potential $\varphi$}{%
      \varphi=\tfrac{W}{Q}\;;\quad
      \varphi_a-\varphi_b=U_{ab}=\int_a^b\!\vec{E}\,\mathrm{d}\vec{s}}\\[0.25em]%
    \formelblock{Wirbelfrei (Elektrostatik)}{%
      \oint\vec{E}\,\mathrm{d}\vec{s}=0
      \;\Leftrightarrow\;\operatorname{rot}\vec{E}=0}\\[0.25em]%
    \formelblock{Potentialgradient}{%
      \vec{E}=-\operatorname{grad}\varphi=-\nabla\varphi}} &
  \beispielblock{$\varphi$ im Homogenfeld}{%
    \begin{aligned}
      &\alpha=90°\Rightarrow W=0\\
      &\Rightarrow\text{Äquipotentialflächen}\\
      &\quad\perp\text{ Feldlinien}\\[0.15em]
      &W_{ab}=Q(\varphi_a-\varphi_b)
    \end{aligned}} \\[0.3em]
  \formelblock{Gradient (kartesisch)}{%
    \vec{E}=-\nabla\varphi=-\!\left(\tfrac{\partial\varphi}{\partial x}\vec{e}_x
      +\tfrac{\partial\varphi}{\partial y}\vec{e}_y
      +\tfrac{\partial\varphi}{\partial z}\vec{e}_z\right)} &
  \formelblock{Divergenz $\vec{D}$ (kartesisch)}{%
    \operatorname{div}\vec{D}=\tfrac{\partial D_x}{\partial x}
      +\tfrac{\partial D_y}{\partial y}+\tfrac{\partial D_z}{\partial z}=\nabla\cdot\vec{D}=\rho} &
  \beispielblock{Zusammenfassung Elektrostatik}{%
    \begin{aligned}
      \operatorname{rot}\vec{E}&=0\quad\text{(wirbelfrei)}\\
      \vec{E}&=-\operatorname{grad}\varphi\\
      \operatorname{div}\vec{D}&=\rho\\
      \vec{D}&=\varepsilon\vec{E}
    \end{aligned}} \\[0.3em]
\end{formelreihe}
\bereichende

% 2.4  Kapazität
\bereich[cKap2]{2.4 Kapazität}
\begin{formelreihe}
  \parbox[t]{\linewidth}{%
    \formelblock{Definition}{%
      C=\tfrac{Q}{U}\;,\quad[C]=1\,\tfrac{\text{As}}{\text{V}}=1\,\text{F}}\\[0.25em]%
    \formelblock{Längenbez. Kapazität}{%
      C'=\tfrac{C}{l}=\tfrac{\lambda}{U}\quad\left[\tfrac{\text{As}}{\text{m}^2}\right]}\\[0.25em]%
    \formelblock{Allgemein}{%
      C=\tfrac{Q}{U}=\tfrac{\oint\vec{D}\,\mathrm{d}\vec{A}}{\int\vec{E}\,\mathrm{d}\vec{s}}}} &
  \parbox[t]{\linewidth}{%
    \formelblock{Plattenkondensator}{%
      C=\varepsilon\cdot\tfrac{A}{d}}\\[0.25em]%
    \formelblock{Zylinderkondensator}{%
      \begin{aligned}
        C&=\tfrac{2\pi\varepsilon l}{\ln(r_a/r_i)}\;,\quad
        C'=\tfrac{2\pi\varepsilon}{\ln(r_a/r_i)}\\[0.15em]
        E(r)&=\tfrac{\lambda}{2\pi\varepsilon r}=\tfrac{U}{r\ln(r_a/r_i)}\\[0.15em]
        \varphi(r)&=\tfrac{\lambda}{2\pi\varepsilon}\ln\tfrac{r_a}{r}
      \end{aligned}}} &
  \parbox[t]{\linewidth}{%
    \formelblock{Doppelleitung (Radien $r$, Abstand $a$)}{%
      C=\tfrac{\pi\varepsilon l}{\ln(a/r)}\;,\quad C'=\tfrac{\pi\varepsilon}{\ln(a/r)}}\\[0.25em]%
    \formelblock{Einzelleitung über Erde (Höhe $h$)}{%
      C=\tfrac{2\pi\varepsilon l}{\ln(2h/r)}}\\[0.25em]%
    \infoblock{$C=f(\text{Geometrie, Material})$\\
      $\neq f(U,I,E,D,\ldots)$}\\[0.25em]%
    \grafik[scale=0.72]{%
      % Plattenkondensator
      \draw[line width=0.7pt,red!70!black] (0,0)--(0,0.8);
      \draw[line width=0.7pt,blue!70!black] (0.9,0)--(0.9,0.8);
      \foreach \y in {0.16,0.4,0.64}{\draw[fld] (0.06,\y)--(0.84,\y);}
      \node[font=\tiny] at (0.45,-0.18){Platten};
      % Zylinderkondensator
      \draw[red!70!black] (2.05,0.4) circle(0.42);
      \draw[blue!70!black] (2.05,0.4) circle(0.13);
      \foreach \a in {0,45,...,315}{\draw[fld] ($(2.05,0.4)+(\a:0.15)$)--($(2.05,0.4)+(\a:0.4)$);}
      \node[font=\tiny] at (2.05,-0.18){Zylinder};}} \\[0.3em]
\end{formelreihe}
\bereichende

% 2.5  Energie
\bereich[cKap2]{2.5 Energie im elektrostatischen Feld}
\begin{formelreihe}
  \parbox[t]{\linewidth}{%
    \formelblock{Energiedichte}{%
      w_e=\tfrac{1}{2}\varepsilon E^2=\tfrac{1}{2}D\cdot E=\tfrac{D^2}{2\varepsilon}}\\[0.25em]%
    \formelblock{Energie im Volumen}{%
      W_e=\int w_e\,\mathrm{d}V}} &
  \parbox[t]{\linewidth}{%
    \formelblock{Rotation $\vec{E}$ (Wirbel, kartesisch)}{%
      \operatorname{rot}\vec{E}=\nabla\times\vec{E}=
      \begin{pmatrix}
        \partial_y E_z-\partial_z E_y\\
        \partial_z E_x-\partial_x E_z\\
        \partial_x E_y-\partial_y E_x
      \end{pmatrix}}\\[0.25em]%
    \formelblock{Elektrostatik: wirbelfrei}{%
      W_e=Q\!\oint_{K_A}\!\vec{E}\,\mathrm{d}\vec{s}=0}} &
  \beispielblock{Plattenkond. $A,d$}{%
    \begin{aligned}
      w_e&=\tfrac{1}{2}\varepsilon E^2\\
      W_e&=w_e\cdot A\cdot d\\
         &=\tfrac{1}{2}CU^2
    \end{aligned}} \\[0.3em]
\end{formelreihe}
\bereichende

% 2.6  Grenzflächen
\bereich[cKap2]{2.6 Grenzflächen im elektrostatischen Feld}
\begin{formelreihe}
  \parbox[t]{\linewidth}{%
    \formelblock{Tangentialkomp. stetig (aus $\oint\vec{E}\,\mathrm{d}\vec{s}=0$)}{%
      E_{t1}=E_{t2}}\\[0.25em]%
    \formelblock{Metall ($\kappa\to\infty$)}{%
      E_t=0}} &
  \parbox[t]{\linewidth}{%
    \formelblock{Normalkomp. stetig (aus $\oint\vec{D}\,\mathrm{d}\vec{A}=0$)}{%
      D_{n1}=D_{n2}}\\[0.25em]%
    \formelblock{Brechung der Feldlinien}{%
      \tfrac{\tan\alpha_1}{\tan\alpha_2}=\tfrac{\varepsilon_1}{\varepsilon_2}}} &
  \beispielblock{Grenzfläche $\varepsilon_1|\varepsilon_2$}{%
    \begin{aligned}
      \varepsilon_1 E_{n1}&=\varepsilon_2 E_{n2}\\
      E_{t1}&=E_{t2}\\[0.1em]
      &\Rightarrow\text{Brechung}
    \end{aligned}} \\[0.3em]
\end{formelreihe}
\bereichende

% 2.7  Vektoroperatoren
\bereich[cKap2]{2.7 Vektoroperatoren — Kartesische Koordinaten $(x,y,z)$}
\begin{formelreihe}
  \parbox[t]{\linewidth}{%
    \formelblock{Linien-/Flächen-/Volumenelement}{%
      \begin{aligned}
        \mathrm{d}\vec{s}&=\mathrm{d}x\,\vec{e}_x+\mathrm{d}y\,\vec{e}_y+\mathrm{d}z\,\vec{e}_z\\
        \mathrm{d}V&=\mathrm{d}x\,\mathrm{d}y\,\mathrm{d}z
      \end{aligned}}\\[0.25em]%
    \formelblock{Nabla-Operator}{%
      \nabla=\tfrac{\partial}{\partial x}\vec{e}_x
        +\tfrac{\partial}{\partial y}\vec{e}_y
        +\tfrac{\partial}{\partial z}\vec{e}_z}} &
  \parbox[t]{\linewidth}{%
    \formelblock{Gradient (Skalarfeld $\varphi$)}{%
      \operatorname{grad}\varphi=\nabla\varphi=
        \tfrac{\partial\varphi}{\partial x}\vec{e}_x
        +\tfrac{\partial\varphi}{\partial y}\vec{e}_y
        +\tfrac{\partial\varphi}{\partial z}\vec{e}_z}\\[0.25em]%
    \formelblock{Divergenz}{%
      \operatorname{div}\vec{D}=\nabla\cdot\vec{D}=
        \tfrac{\partial D_x}{\partial x}+\tfrac{\partial D_y}{\partial y}
        +\tfrac{\partial D_z}{\partial z}}\\[0.25em]%
    \formelblock{Laplace}{%
      \Delta=\tfrac{\partial^2}{\partial x^2}
        +\tfrac{\partial^2}{\partial y^2}+\tfrac{\partial^2}{\partial z^2}}} &
  \beispielblock{Bsp: $\varphi=\varphi_0{-}ax$}{%
    \begin{aligned}
      \vec{E}&=-\nabla\varphi=a\,\vec{e}_x\\
      &\Rightarrow\text{homogenes Feld}\\
      &\quad\text{in }x\text{-Richtung}
    \end{aligned}} \\[0.3em]
\end{formelreihe}
\bereichende

\bereich[cKap2]{2.7 Vektoroperatoren — Zylinderkoordinaten $(r,\alpha,z)$}
\begin{formelreihe}
  \parbox[t]{\linewidth}{%
    \formelblock{Linien-/Volumenelement}{%
      \begin{aligned}
        \mathrm{d}\vec{s}&=\mathrm{d}r\,\vec{e}_r+r\,\mathrm{d}\alpha\,\vec{e}_\alpha
          +\mathrm{d}z\,\vec{e}_z\\
        \mathrm{d}V&=r\,\mathrm{d}r\,\mathrm{d}\alpha\,\mathrm{d}z
      \end{aligned}}\\[0.25em]%
    \formelblock{Gradient}{%
      \nabla\varphi=\tfrac{\partial\varphi}{\partial r}\vec{e}_r
        +\tfrac{1}{r}\tfrac{\partial\varphi}{\partial\alpha}\vec{e}_\alpha
        +\tfrac{\partial\varphi}{\partial z}\vec{e}_z}} &
  \parbox[t]{\linewidth}{%
    \formelblock{Divergenz}{%
      \operatorname{div}\vec{D}=
        \tfrac{1}{r}\tfrac{\partial(r D_r)}{\partial r}
        +\tfrac{1}{r}\tfrac{\partial D_\alpha}{\partial\alpha}
        +\tfrac{\partial D_z}{\partial z}}\\[0.25em]%
    \formelblock{Laplace}{%
      \Delta=\tfrac{1}{r}\tfrac{\partial}{\partial r}\!\left(r\tfrac{\partial}{\partial r}\right)
        +\tfrac{1}{r^2}\tfrac{\partial^2}{\partial\alpha^2}
        +\tfrac{\partial^2}{\partial z^2}}} &
  \beispielblock{Zylinderkondensator}{%
    \begin{aligned}
      D(r)&=\tfrac{\lambda}{2\pi r}\\
      E(r)&=\tfrac{\lambda}{2\pi\varepsilon r}\\
      \varphi(r)&=\tfrac{\lambda}{2\pi\varepsilon}\ln\tfrac{r_a}{r}
    \end{aligned}} \\[0.3em]
\end{formelreihe}
\bereichende

\bereich[cKap2]{2.7 Vektoroperatoren — Kugelkoordinaten $(r,\vartheta,\alpha)$}
\begin{formelreihe}
  \parbox[t]{\linewidth}{%
    \formelblock{Volumenelement}{%
      \mathrm{d}V=r^2\sin\vartheta\,\mathrm{d}r\,\mathrm{d}\vartheta\,\mathrm{d}\alpha}\\[0.25em]%
    \formelblock{Gradient}{%
      \nabla\varphi=\tfrac{\partial\varphi}{\partial r}\vec{e}_r
        +\tfrac{1}{r}\tfrac{\partial\varphi}{\partial\vartheta}\vec{e}_\vartheta
        +\tfrac{1}{r\sin\vartheta}\tfrac{\partial\varphi}{\partial\alpha}\vec{e}_\alpha}} &
  \parbox[t]{\linewidth}{%
    \formelblock{Divergenz}{%
      \operatorname{div}\vec{D}=
        \tfrac{1}{r^2}\tfrac{\partial(r^2 D_r)}{\partial r}
        +\tfrac{1}{r\sin\vartheta}\tfrac{\partial(\sin\vartheta\, D_\vartheta)}{\partial\vartheta}
        +\tfrac{1}{r\sin\vartheta}\tfrac{\partial D_\alpha}{\partial\alpha}}\\[0.25em]%
    \formelblock{Laplace}{%
      \Delta=\tfrac{1}{r^2}\tfrac{\partial}{\partial r}\!\left(r^2\tfrac{\partial}{\partial r}\right)
        +\tfrac{1}{r^2\sin\vartheta}\tfrac{\partial}{\partial\vartheta}\!\left(\sin\vartheta
        \tfrac{\partial}{\partial\vartheta}\right)
        +\tfrac{1}{r^2\sin^2\vartheta}\tfrac{\partial^2}{\partial\alpha^2}}} &
  \beispielblock{Punktladung $Q$}{%
    \begin{aligned}
      \vec{D}&=\tfrac{Q}{4\pi r^2}\vec{e}_r\\
      \vec{E}&=\tfrac{Q}{4\pi\varepsilon r^2}\vec{e}_r\\
      \varphi&=\tfrac{Q}{4\pi\varepsilon r}
    \end{aligned}} \\[0.3em]
\end{formelreihe}
\bereichende

% ---- EDy3: Magnetisches Feld ----------------------------

% 3.1  Magnetische Feldstärke
\bereich[cKap3]{3.1 Magnetische Feldstärke $\vec{H}$}
\begin{formelreihe}
  \parbox[t]{\linewidth}{%
    \formelblock{Magnetische Feldstärke}{%
      H=|\vec{H}|=\tfrac{I}{2\pi r}=\tfrac{\text{Stromstärke}}{\text{Feldlinienlänge}}}\\[0.25em]%
    \formelblock{Einheit}{%
      [H]=1\,\tfrac{\mathrm{A}}{\mathrm{m}}}\\[0.25em]%
    \infoblock{$\vec{H}\perp I$;\; $\vec{H}\perp\vec{r}$\quad (Rechtsschraube)}} &
  \parbox[t]{\linewidth}{%
    \formelblock{Proportionalitäten}{%
      H\sim I\;,\quad H\sim\tfrac{1}{r}}\\[0.25em]%
    \infoblock{analog zu $E=U/d$:\\
      $H=U_{\mathrm{mag}}/l_{\mathrm{Feldlinie}}$}} &
  \beispielblock{Langer gerader Leiter}{%
    \begin{aligned}
      H(r)&=\tfrac{I}{2\pi r}\\[0.1em]
      [H]&=1\,\tfrac{\mathrm{A}}{\mathrm{m}}
    \end{aligned}} \\[0.3em]
\end{formelreihe}
\bereichende

% 3.2  Magnetische Flussdichte
\bereich[cKap3]{3.2 Magnetische Flussdichte $\vec{B}$}
\begin{formelreihe}
  \parbox[t]{\linewidth}{%
    \formelblock{Flussdichte und Fluss}{%
      \vec{B}=\tfrac{\mathrm{d}\Phi}{\mathrm{d}\vec{A}}\;,\qquad
      \Phi=\int\!\vec{B}\,\mathrm{d}\vec{A}}\\[0.25em]%
    \formelblock{Einheit}{%
      [B]=1\,\tfrac{\mathrm{Vs}}{\mathrm{m}^2}=1\,\mathrm{T}}} &
  \parbox[t]{\linewidth}{%
    \formelblock{Materialgesetz}{%
      \vec{B}=\mu\cdot\vec{H}=\mu_0\cdot\mu_r\cdot\vec{H}}\\[0.25em]%
    \formelblock{Permeabilität Vakuum}{%
      \mu_0=4\pi\cdot10^{-7}\,\tfrac{\mathrm{Vs}}{\mathrm{Am}}
        =1{,}257\cdot10^{-6}\,\tfrac{\mathrm{Vs}}{\mathrm{Am}}}} &
  \beispielblock{analog Elektrostatik}{%
    \begin{aligned}
      \vec{D}&=\varepsilon\vec{E}
        \;\leftrightarrow\;\vec{B}=\mu\vec{H}\\[0.1em]
      I&=\int\vec{J}\,\mathrm{d}\vec{A}
        \;\leftrightarrow\;\Phi=\int\vec{B}\,\mathrm{d}\vec{A}
    \end{aligned}} \\[0.3em]
\end{formelreihe}
\bereichende

% 3.3  Durchflutungsgesetz
\bereich[cKap3]{3.3 Durchflutungsgesetz}
\begin{formelreihe}
  \parbox[t]{\linewidth}{%
    \formelblock{Durchflutungsgesetz}{%
      I=\int_A\!\vec{J}\,\mathrm{d}\vec{A}
        =\oint_{K_A}\!\vec{H}\,\mathrm{d}\vec{s}=V}\\[0.25em]%
    \formelblock{Magn.\ Spannung (homogen)}{%
      V=H\cdot l}} &
  \parbox[t]{\linewidth}{%
    \formelblock{Weg umschließt Leiter}{%
      V_0=\oint\vec{H}\,\mathrm{d}\vec{s}=I}\\[0.25em]%
    \formelblock{Weg umschließt Leiter \emph{nicht}}{%
      \oint\vec{H}\,\mathrm{d}\vec{s}=0}} &
  \beispielblock{$\vec{H}$: Wirbelfeld}{%
    \begin{aligned}
      &\oint\vec{H}\,\mathrm{d}\vec{s}=I\neq0\\[0.1em]
      &\Rightarrow\text{nicht wirbelfrei}\\[0.1em]
      &\text{vgl.\ Elektrostatik:}\\
      &\oint\vec{E}\,\mathrm{d}\vec{s}=0
    \end{aligned}} \\[0.3em]
\end{formelreihe}
\bereichende

% 3.4.1  Bewegungsinduktion
\bereich[cKap3]{3.4.1 Bewegungsinduktion (Leiter bewegt, $\vec{B}$ ruhend)}
\begin{formelreihe}
  \parbox[t]{\linewidth}{%
    \formelblock{Induziertes Feld (Lorentzkraft)}{%
      \vec{E}_i=\vec{v}\times\vec{B}}\\[0.25em]%
    \formelblock{Induz.\ Spannung (homogen, $\vec{v}\perp\vec{B}\perp\vec{l}$)}{%
      U_i=-B\cdot l\cdot v}} &
  \parbox[t]{\linewidth}{%
    \formelblock{Allgemein (Umlaufintegral)}{%
      U_i=\oint\vec{E}_i\,\mathrm{d}\vec{s}
        =\oint\!\left(\vec{v}\times\vec{B}\right)\mathrm{d}\vec{s}}\\[0.25em]%
    \formelblock{Induktive Spannung an Klemmen}{%
      U_q=-U_i\quad\text{(EMK)}}} &
  \beispielblock{Betrag $E_i$}{%
    \begin{aligned}
      &\vec{v}\perp\vec{B}:\;E_{max}=v\cdot B\\[0.1em]
      &\vec{v}\parallel\vec{B}:\;E_i=0
    \end{aligned}} \\[0.3em]
\end{formelreihe}
\bereichende

% 3.4.2 / 3.4.3  Ruheinduktion & allgemeines Induktionsgesetz
\bereich[cKap3]{3.4.2 Ruheinduktion \& 3.4.3 Allgemeines Induktionsgesetz}
\begin{formelreihe}
  \parbox[t]{\linewidth}{%
    \formelblock{Induktionsgesetz ($N$ Windungen)}{%
      u_q=N\cdot\tfrac{\mathrm{d}\Phi}{\mathrm{d}t}}\\[0.25em]%
    \infoblock{\textbf{Lenzsche Regel:} Induzierter Strom\\
      wirkt der Flussänderung entgegen.}} &
  \parbox[t]{\linewidth}{%
    \formelblock{2.\ Maxwellsche Gleichung}{%
      u_i=\oint\vec{E}_i\,\mathrm{d}\vec{s}
        =-\tfrac{\mathrm{d}}{\mathrm{d}t}\int_A\!\vec{B}\,\mathrm{d}\vec{A}}\\[0.25em]%
    \formelblock{Produktregel ($\Phi=B(t)\cdot A(t)$)}{%
      -u_i=\underbrace{\tfrac{\mathrm{d}B}{\mathrm{d}t}\cdot A}_{\text{Ruhe}}
        +\underbrace{B\cdot\tfrac{\mathrm{d}A}{\mathrm{d}t}}_{\text{Bewegung}}}} &
  \beispielblock{Anteile}{%
    \begin{aligned}
      \tfrac{\mathrm{d}B}{\mathrm{d}t}{\cdot}A
        &:\;\text{Ruheinduktion}\\
        &\quad\text{(Trafo)}\\[0.1em]
      B{\cdot}\tfrac{\mathrm{d}A}{\mathrm{d}t}
        &:\;\text{Bewegungsinduktion}\\
        &\quad\text{(Generator)}
    \end{aligned}} \\[0.3em]
\end{formelreihe}
\bereichende

% 3.5  Induktivität
\bereich[cKap3]{3.5 Induktivität}
\begin{formelreihe}
  \parbox[t]{\linewidth}{%
    \formelblock{Selbstinduktivität}{%
      L=\tfrac{\psi_m}{I}=\tfrac{N\cdot\Phi}{I}\;,\quad
      [L]=1\,\tfrac{\mathrm{Vs}}{\mathrm{A}}=1\,\mathrm{H}}\\[0.25em]%
    \formelblock{Spulenspannung}{%
      u_L(t)=L\cdot\tfrac{\mathrm{d}i}{\mathrm{d}t}=-u_i(t)}} &
  \parbox[t]{\linewidth}{%
    \formelblock{Aus Geometrie \& Material}{%
      L=\mu\cdot N^2\cdot\tfrac{A}{l}}\\[0.25em]%
    \formelblock{Doppelleitung (exakt / $a\gg r$)}{%
      \begin{aligned}
        L&=\tfrac{\mu_0 l}{\pi}\!\left(\ln\tfrac{a-r}{r}+0{,}25\right)\\[0.1em]
        L&\approx\tfrac{\mu_0 l}{\pi}\cdot\ln\tfrac{a}{r}
      \end{aligned}}} &
  \beispielblock{Magn.\ Kreis}{%
    \begin{aligned}
      \Phi&=\tfrac{\Theta}{R_m}\;,\;\Theta=I N\\[0.1em]
      \Lambda_m&=\tfrac{1}{R_m}=\tfrac{\mu A}{l}\\[0.1em]
      L&=\Lambda_m\cdot N^2
    \end{aligned}} \\[0.3em]
\end{formelreihe}
\bereichende

% 3.6  Energie im magnetischen Feld
\bereich[cKap3]{3.6 Energie im magnetischen Feld}
\begin{formelreihe}
  \parbox[t]{\linewidth}{%
    \formelblock{Energie (Spule)}{%
      W_m=\tfrac{1}{2}L I^2=\tfrac{1}{2}I\psi_m
        =\tfrac{1}{2}\Theta\cdot\Phi}} &
  \parbox[t]{\linewidth}{%
    \formelblock{Energiedichte}{%
      w_m=\tfrac{1}{2}\mu H^2=\tfrac{1}{2}B\cdot H
        =\tfrac{B^2}{2\mu}}} &
  \beispielblock{analog Elektrostatik}{%
    \begin{aligned}
      W_e&=\tfrac{1}{2}CU^2\;\leftrightarrow\;W_m=\tfrac{1}{2}LI^2\\[0.1em]
      w_e&=\tfrac{1}{2}\varepsilon E^2\;\leftrightarrow\;w_m=\tfrac{1}{2}\mu H^2
    \end{aligned}} \\[0.3em]
\end{formelreihe}
\bereichende

% 3.7  Grenzflächen im magnetischen Feld
\bereich[cKap3]{3.7 Grenzflächen im magnetischen Feld}
\begin{formelreihe}
  \parbox[t]{\linewidth}{%
    \formelblock{Tangentialkomp.\ stetig (aus $\oint\vec{H}\,\mathrm{d}\vec{s}=0$)}{%
      H_{t1}=H_{t2}}\\[0.25em]%
    \formelblock{Brechung der Feldlinien}{%
      \tfrac{\tan\alpha_1}{\tan\alpha_2}=\tfrac{\mu_1}{\mu_2}}} &
  \parbox[t]{\linewidth}{%
    \formelblock{Normalkomp.\ stetig (aus $\oint\vec{B}\,\mathrm{d}\vec{A}=0$)}{%
      B_{n1}=B_{n2}}\\[0.25em]%
    \formelblock{Grenzfl.\ zu Metall ($\kappa\to\infty$, zeitl.\ veränd.)}{%
      H_n=0}} &
  \beispielblock{Zusammenfassung Magn.\ Feld}{%
    \begin{aligned}
      \oint\vec{H}\,\mathrm{d}\vec{s}&=I\;\text{(Wirbelfeld)}\\
      \oint\vec{B}\,\mathrm{d}\vec{A}&=0\;\text{(quellenfrei)}\\
      \vec{B}&=\mu\vec{H}
    \end{aligned}} \\[0.3em]
\end{formelreihe}
\bereichende

% ---- EDy4: Stationäres Strömungsfeld --------------------

% 4.1  Stromdichte
\bereich[cKap4]{4.1 Stromdichte $\vec{J}$}
\begin{formelreihe}
  \parbox[t]{\linewidth}{%
    \formelblock{Stromdichte (allgemein)}{%
      I=\int_A\!\vec{J}\,\mathrm{d}\vec{A}}\\[0.25em]%
    \formelblock{Homogenes Feld}{%
      J=\tfrac{I}{A}\;,\qquad [J]=1\,\tfrac{\mathrm{A}}{\mathrm{m}^2}}} &
  \parbox[t]{\linewidth}{%
    \formelblock{Stromdichtevektor}{%
      \vec{J}\parallel\vec{E}\;,\quad
      \vec{J}\parallel\vec{S}_{\mathrm{Feldl.}}}\\[0.25em]%
    \infoblock{Strömungsfeld quellenfrei:\\
      $\oint\vec{J}\,\mathrm{d}\vec{A}=0$}} &
  \beispielblock{Draht, $d=2\,\text{mm}$, $I=1\,\text{A}$}{%
    \begin{aligned}
      A&=\pi r^2=\pi(1\,\text{mm})^2\\[0.1em]
      J&=\tfrac{1\,\text{A}}{\pi\cdot10^{-6}\,\text{m}^2}\\[0.1em]
      &\approx318\,\tfrac{\text{kA}}{\text{m}^2}
    \end{aligned}} \\[0.3em]
\end{formelreihe}
\bereichende

% 4.2  Elektrische Leitfähigkeit / Lok. Ohmsches Gesetz
\bereich[cKap4]{4.2 Elektrische Leitfähigkeit — Lokales Ohmsches Gesetz}
\begin{formelreihe}
  \parbox[t]{\linewidth}{%
    \formelblock{Lokales Ohmsches Gesetz}{%
      \vec{J}=\kappa\cdot\vec{E}}\\[0.25em]%
    \formelblock{Spez.\ Leitfähigkeit}{%
      [\kappa]=1\,\tfrac{\mathrm{S}}{\mathrm{m}}
        =1\,\tfrac{\mathrm{A}}{\mathrm{V\,m}}}} &
  \parbox[t]{\linewidth}{%
    \formelblock{Spez.\ Widerstand}{%
      \varrho=\tfrac{1}{\kappa}\;,\qquad
      [\varrho]=1\,\Omega\mathrm{m}}\\[0.25em]%
    \infoblock{Kupfer: $\kappa_\mathrm{Cu}=58\cdot10^6\,\tfrac{\mathrm{S}}{\mathrm{m}}$\\
      Isolator: $\kappa\to0$, Metall: $\kappa\gg1$}} &
  \beispielblock{Analogie}{%
    \begin{aligned}
      \vec{J}&=\kappa\,\vec{E}\\[0.1em]
      \vec{D}&=\varepsilon\,\vec{E}\\[0.1em]
      \Rightarrow\;\kappa&\leftrightarrow\varepsilon
    \end{aligned}} \\[0.3em]
\end{formelreihe}
\bereichende

% 4.3  Kontinuitätsgleichung
\bereich[cKap4]{4.3 Kontinuitätsgleichung (Stationär)}
\begin{formelreihe}
  \parbox[t]{\linewidth}{%
    \formelblock{Quellfreiheit der Stromdichte}{%
      \oint_{A}\!\vec{J}\,\mathrm{d}\vec{A}=0}\\[0.25em]%
    \formelblock{Differenzielle Form}{%
      \nabla\cdot\vec{J}=0}} &
  \parbox[t]{\linewidth}{%
    \formelblock{Knotenpunktsatz (Kirchhoff)}{%
      \textstyle\sum_k I_k=0}\\[0.25em]%
    \infoblock{Stationär: keine freien Ladungen\\
      im Leiterinneren ($\varrho_\mathrm{frei}=0$)}} &
  \beispielblock{Zwei Leiter, $A_1\neq A_2$}{%
    \begin{aligned}
      I_1&=I_2=I\\[0.1em]
      J_1 A_1&=J_2 A_2\\[0.1em]
      \Rightarrow J&\sim\tfrac{1}{A}
    \end{aligned}} \\[0.3em]
\end{formelreihe}
\bereichende

% 4.4  Widerstand — Globales Ohmsches Gesetz
\bereich[cKap4]{4.4 Widerstand — Globales Ohmsches Gesetz}
\begin{formelreihe}
  \parbox[t]{\linewidth}{%
    \formelblock{Ohmsches Gesetz (integral)}{%
      U=R\cdot I\;,\qquad [R]=1\,\Omega}\\[0.25em]%
    \formelblock{Widerstand aus Geometrie}{%
      R=\varrho\cdot\tfrac{l}{A}=\tfrac{l}{\kappa\cdot A}}} &
  \parbox[t]{\linewidth}{%
    \formelblock{Leitwert}{%
      G=\tfrac{1}{R}=\kappa\cdot\tfrac{A}{l}\;,\quad[G]=1\,\mathrm{S}}\\[0.25em]%
    \formelblock{Herleitung ($\vec{J}=\kappa\vec{E}$, homogen)}{%
      U=E\cdot l\;,\;I=J\cdot A=\kappa E A}} &
  \beispielblock{Zylinder, $l$, $A$, $\kappa$}{%
    \begin{aligned}
      R&=\tfrac{l}{\kappa A}\\[0.1em]
      U&=\tfrac{I\,l}{\kappa A}\\[0.1em]
      E&=\tfrac{U}{l}=\tfrac{J}{\kappa}
    \end{aligned}} \\[0.3em]
\end{formelreihe}
\bereichende

% 4.5  Verlustleistung (Joulesche Wärme)
\bereich[cKap4]{4.5 Verlustleistung — Joulesche Wärme}
\begin{formelreihe}
  \parbox[t]{\linewidth}{%
    \formelblock{Leistung (global)}{%
      P=U\cdot I=R\cdot I^2=\tfrac{U^2}{R}}\\[0.25em]%
    \formelblock{Leistungsdichte (lokal)}{%
      p=\kappa\cdot E^2=\vec{J}\cdot\vec{E}=\tfrac{J^2}{\kappa}}} &
  \parbox[t]{\linewidth}{%
    \formelblock{Gesamtleistung aus Dichte}{%
      P=\int_V\!p\,\mathrm{d}V
        =\int_V\!\kappa E^2\,\mathrm{d}V}\\[0.25em]%
    \infoblock{Elektrische Energie → Wärme\\
      (Joulsche Verluste, irreversibel)}} &
  \beispielblock{Zylinder, $\kappa$, $E$, $V=A\cdot l$}{%
    \begin{aligned}
      p&=\kappa E^2\\[0.1em]
      P&=\kappa E^2\cdot A l\\[0.1em]
      &=\tfrac{U^2}{R}
    \end{aligned}} \\[0.3em]
\end{formelreihe}
\bereichende

% 4.6  Grenzflächen im Strömungsfeld
\bereich[cKap4]{4.6 Grenzflächen im Strömungsfeld}
\begin{formelreihe}
  \parbox[t]{\linewidth}{%
    \formelblock{Normalkomp.\ stetig ($\oint\vec{J}\,\mathrm{d}\vec{A}=0$)}{%
      J_{n1}=J_{n2}}\\[0.25em]%
    \formelblock{Tangentialkomp.\ stetig ($\oint\vec{E}\,\mathrm{d}\vec{s}=0$)}{%
      E_{t1}=E_{t2}}} &
  \parbox[t]{\linewidth}{%
    \formelblock{Brechungsgesetz}{%
      \tfrac{\tan\alpha_1}{\tan\alpha_2}=\tfrac{\kappa_1}{\kappa_2}}\\[0.25em]%
    \infoblock{Analog zum Magnetfeld:\\
      $\tfrac{\tan\alpha_1}{\tan\alpha_2}=\tfrac{\mu_1}{\mu_2}$\quad
      ($J_{n}\leftrightarrow B_{n}$)}} &
  \beispielblock{Grenzfl.\ $\kappa_1\gg\kappa_2$}{%
    \begin{aligned}
      &\kappa_1\gg\kappa_2:\\[0.1em]
      &\Rightarrow\alpha_1\gg\alpha_2\\[0.1em]
      &\Rightarrow\vec{J}\text{ senkr.\ in Med.\ 2}
    \end{aligned}} \\[0.3em]
\end{formelreihe}
\bereichende

% 4.7  Analogie Strömungsfeld ↔ Elektrostatik
\bereich[cKap4]{4.7 Analogie: Strömungsfeld $\leftrightarrow$ Elektrostatik}
\begin{formelreihe}
  \parbox[t]{\linewidth}{%
    \formelblock{Feldgrößen-Analogie}{%
      \vec{J}=\kappa\vec{E}
        \;\leftrightarrow\;
      \vec{D}=\varepsilon\vec{E}}\\[0.25em]%
    \formelblock{Quell-/Wirbelfreiheit}{%
      \oint\vec{J}\,\mathrm{d}\vec{A}=0
        \;\leftrightarrow\;
      \oint\vec{D}\,\mathrm{d}\vec{A}=0}} &
  \parbox[t]{\linewidth}{%
    \formelblock{Schaltgrößen-Analogie}{%
      G\leftrightarrow C\;,\quad R\leftrightarrow\tfrac{1}{C}}\\[0.25em]%
    \formelblock{Zusammenhang $R$ und $C$}{%
      R\cdot C=\tfrac{\varepsilon}{\kappa}}} &
  \beispielblock{Plattenkondens.\ $A$, $d$}{%
    \begin{aligned}
      C&=\varepsilon\tfrac{A}{d}\\[0.1em]
      G&=\kappa\tfrac{A}{d}\\[0.1em]
      R\cdot C&=\tfrac{\varepsilon}{\kappa}
    \end{aligned}} \\[0.3em]
\end{formelreihe}
\bereichende

%% ---- EDy5: Zeitveränderliches Feld (Wechselstrom) --------

% 5.1  Maxwellsche Gleichungen — klausurrelevant
\bereich[cKap5]{5.1 Maxwellsche Gleichungen \; (klausurrelevant!)}
\begin{formelreihe}
  \parbox[t]{\linewidth}{%
    \merkblock[cKap5]{Differenzielle Form}{%
      \begin{aligned}
        \nabla\times\vec H &= \tfrac{\mathrm{d}\vec D}{\mathrm{d}t}+\vec J\\[0.15em]
        \nabla\times\vec E &= -\tfrac{\mathrm{d}\vec B}{\mathrm{d}t}\\[0.15em]
        \nabla\cdot\vec D &= \varrho\\[0.15em]
        \nabla\cdot\vec B &= 0
      \end{aligned}}\\[0.25em]%
    \formelblock{Kontinuitätsgleichung}{%
      \nabla\cdot\vec J=-\tfrac{\mathrm{d}\varrho}{\mathrm{d}t}}} &
  \parbox[t]{\linewidth}{%
    \merkblock[cKap5]{Integrale Form}{%
      \begin{aligned}
        \oint\vec H\,\mathrm{d}\vec s &= \int_A\!\tfrac{\mathrm{d}\vec D}{\mathrm{d}t}\,\mathrm{d}\vec A+I\\[0.15em]
        \oint\vec E\,\mathrm{d}\vec s &= -\int_A\!\tfrac{\mathrm{d}\vec B}{\mathrm{d}t}\,\mathrm{d}\vec A\\[0.15em]
        \oint_A\!\vec D\,\mathrm{d}\vec A &= Q\\[0.15em]
        \oint_A\!\vec B\,\mathrm{d}\vec A &= 0
      \end{aligned}}} &
  \parbox[t]{\linewidth}{%
    \formelblock{Bedeutung (von oben nach unten)}{%
      \begin{aligned}
        &\text{1. Durchflutungsgesetz}\\
        &\text{2. Induktionsgesetz}\\
        &\text{3. Ladungen sind Quellen}\\
        &\text{4. } \vec B \text{ quellenfrei}
      \end{aligned}}\\[0.25em]%
    \formelblock{Materialgleichungen}{%
      \vec D=\varepsilon_0\varepsilon_r\vec E\;,\quad
      \vec B=\mu_0\mu_r\vec H}\\[0.2em]%
    \formelblock{Lokales Ohmsches Gesetz}{%
      \vec J=\kappa\cdot\vec E}} \\[0.3em]
\end{formelreihe}
\bereichende

% 5.2  Ausbreitung im Leiter — Fortpflanzungskonstante
\bereich[cKap5]{5.2 Ausbreitung im Leiter — Fortpflanzungskonstante}
\begin{formelreihe}
  \parbox[t]{\linewidth}{%
    \formelblock{Komplexer Ansatz (Wechselstrom)}{%
      \underline{\vec E}=\vec E(\vec r)\,e^{j\varphi}\cdot e^{j\omega t}
        \quad(\text{ebenso } \underline{\vec H},\underline{\vec J})}\\[0.25em]%
    \infoblock{Quaderleiter, Strom in $z$, $x\!\ge\!0$ in den Leiter:\\
      $\underline{\vec J}=\underline J_z(x)\vec e_z$,\;
      $\underline{\vec H}=\underline H_y(x)\vec e_y$,\;
      $\underline{\vec E}=\underline E_z(x)\vec e_z$.\\
      Kein Verschiebungsstrom im Leiter ($\kappa\!\gg\!\omega\varepsilon$).}} &
  \parbox[t]{\linewidth}{%
    \formelblock{(1) Durchflutungsgesetz}{%
      \underline J_z=\tfrac{\mathrm{d}\underline H_y}{\mathrm{d}x}
        =\kappa\,\underline E_z}\\[0.25em]%
    \formelblock{(2) Induktionsgesetz}{%
      \tfrac{\mathrm{d}\underline E_z}{\mathrm{d}x}=j\omega\mu\,\underline H_y}\\[0.25em]%
    \formelblock{Kombiniert (DGL 2.\ Ordnung)}{%
      \tfrac{\mathrm{d}^2\underline E_z}{\mathrm{d}x^2}=j\omega\mu\kappa\,\underline E_z}} &
  \parbox[t]{\linewidth}{%
    \formelblock{Lösung (abklingend in $+x$)}{%
      \underline E_z(x)=\underline E_1\,e^{-\underline\gamma x}}\\[0.25em]%
    \merkblock[cKap5]{Fortpflanzungskonstante}{%
      \underline\gamma=\sqrt{j\omega\mu\kappa}
        =\sqrt{\tfrac{\omega\mu\kappa}{2}}\,(1+j)=\alpha+j\beta}\\[0.2em]%
    \infoblock{$\alpha=\beta=\sqrt{\tfrac{\omega\mu\kappa}{2}}=\tfrac{1}{\delta}$\\
      $\alpha$: Amplitudenabnahme,\; $\beta$: Phasendrehung.}} \\[0.3em]
\end{formelreihe}
\bereichende

% 5.3  Skineffekt / Skintiefe
\bereich[cKap5]{5.3 Skineffekt — Skintiefe (Eindringtiefe) $\delta$}
\begin{formelreihe}
  \parbox[t]{\linewidth}{%
    \merkblock[cKap5]{Skintiefe (äquiv.\ Leitschichtdicke)}{%
      \delta=\sqrt{\tfrac{2}{\omega\mu\kappa}}\;,\qquad [\delta]=\mathrm{m}}\\[0.25em]%
    \formelblock{Reale Feldstärkeverteilung}{%
      E_z(x,t)=E_1\,e^{-x/\delta}\cos\!\big(\omega t+\varphi_E-\tfrac{x}{\delta}\big)}\\[0.2em]%
    \infoblock{Amplitude fällt auf $e^{-1}$ nach Tiefe $x=\delta$.\\
      Strom fließt überwiegend in der Randschicht
      ($\to$ \emph{Hauteffekt}).}} &
  \parbox[t]{\linewidth}{%
    \formelblock{$\delta$ ist abhängig von}{%
      \begin{aligned}
        &\omega:\ \text{Frequenz}\ (\delta\!\sim\!1/\!\sqrt{f})\\
        &\mu:\ \text{magn.\ Eigenschaften}\\
        &\kappa:\ \text{Leitfähigkeit}
      \end{aligned}}\\[0.2em]%
    \infoblock{$\to$ hohe $f$, hohe $\kappa$, hohe $\mu$
      $\Rightarrow$ kleine $\delta$.}} &
  \parbox[t]{\linewidth}{%
    \formelblock{Beispielwerte $\delta$ (Kupfer $\kappa_{\mathrm{Cu}}$)}{%
      \begin{array}{@{}l l@{}}
        50\,\text{Hz}: & \delta_{\mathrm{Cu}}\approx9{,}5\,\text{mm}\\
        & \delta_{\mathrm{Al}}\approx12\,\text{mm},\ \delta_{\mathrm{Fe}}\!\approx\!0{,}6\,\text{mm}\\
        100\,\text{MHz}: & \delta_{\mathrm{Cu}}\approx6{,}7\,\mu\mathrm{m}
      \end{array}}\\[0.2em]%
    \infoblock{Eisen trotz kleinem $\kappa$ dünn, da $\mu_r\!\gg\!1$.}} \\[0.3em]
\end{formelreihe}
\bereichende

% 5.4  Wechselstromwiderstand
\bereich[cKap5]{5.4 Wechselstromwiderstand — innere Induktivität}
\begin{formelreihe}
  \parbox[t]{\linewidth}{%
    \formelblock{Gleichstromwiderstand ($b\times d$, Länge $l$)}{%
      R_{=}=\tfrac{l}{\kappa\,b\,d}
        \qquad(\text{Runddraht: } \tfrac{4l}{\pi d^2\kappa})}\\[0.25em]%
    \merkblock[cKap5]{Impedanz (dicker Leiter, $\delta\!\ll\!d$)}{%
      \underline Z=\tfrac{l}{\kappa\,b\,\delta}(1+j)=R_{\approx}+jX_i}\\[0.2em]%
    \infoblock{Real- und Imaginärteil betragsgleich.\\
      $X_i$: \emph{innere Induktivität} des Leiters.}} &
  \parbox[t]{\linewidth}{%
    \formelblock{Wechselstromwiderstand}{%
      R_{\approx}=\tfrac{l}{\kappa\,b\,\delta}
        =\tfrac{l}{b}\cdot R_F}\\[0.25em]%
    \formelblock{Oberflächen-/Flächenwiderstand}{%
      R_F=\tfrac{1}{\kappa\,\delta}=R_{\square}\;,\quad[R_F]=\Omega}\\[0.2em]%
    \merkblock{Äquivalenz}{d=\delta\;\Rightarrow\;R_{\approx}(d\!\to\!\infty)=R_{=}(\text{Dicke }\delta)}} &
  \parbox[t]{\linewidth}{%
    \formelblock{Runddraht — Näherung ($R_{=}=\tfrac{4l}{\pi d^2\kappa}$)}{%
      \begin{array}{@{}l@{\;}l@{}}
        d<2\delta: & R_{\approx}\!\approx\!R_{=}\\[0.15em]
        2\delta\!\le\!d\!<\!4\delta: & R_{\approx}\!\approx\!R_{=}\!\big[1\!+\!(\tfrac{d}{5{,}3\delta})^4\big]\\[0.15em]
        4\delta\!\le\!d\!<\!10\delta: & R_{\approx}\!\approx\!R_{=}\!\big[\tfrac{d}{4\delta}\!+\!0{,}25\big]\\[0.15em]
        10\delta\!\le\!d: & R_{\approx}\!\approx\!R_{=}\tfrac{d}{4\delta}
      \end{array}}\\[0.2em]%
    \infoblock{$\to$ bei hohen $f$ viele dünne Leiter statt
      eines dicken (\emph{HF-Litze}); massive Leiter $>\!20$\,mm
      auch bei 50\,Hz nicht mehr optimal.}} \\[0.3em]
\end{formelreihe}
\bereichende

%% ---- EDy6: Elektromagnetische Wellen --------------------

% 6.1  Wellengleichungen
\bereich[cKap6]{6.1 Wellengleichungen}
\begin{formelreihe}
  \parbox[t]{\linewidth}{%
    \formelblock{Ausgangsgleichungen ($\varepsilon,\mu,\kappa=$const)}{%
      \begin{aligned}
        \nabla\times\vec H&=\kappa\vec E+\varepsilon\tfrac{\partial\vec E}{\partial t}\\[0.1em]
        \nabla\times\vec E&=-\mu\tfrac{\partial\vec H}{\partial t}
      \end{aligned}}\\[0.2em]%
    \infoblock{Herleitung wie beim Skineffekt: pro Feld\-komponente
      eine DGL 2.\ Ordnung ($\mathrm{rot\,rot}=\mathrm{grad\,div}-\Delta$).}} &
  \parbox[t]{\linewidth}{%
    \formelblock{Allg.\ Wellengleichung (Zeitbereich)}{%
      \begin{aligned}
        \Delta\vec H-\kappa\mu\tfrac{\partial\vec H}{\partial t}
          -\varepsilon\mu\tfrac{\partial^2\vec H}{\partial t^2}&=0\\[0.1em]
        \Delta\vec E-\kappa\mu\tfrac{\partial\vec E}{\partial t}
          -\varepsilon\mu\tfrac{\partial^2\vec E}{\partial t^2}&=\mathrm{grad}\tfrac{\varrho}{\varepsilon}
      \end{aligned}}\\[0.2em]%
    \merkblock[cKap6]{Helmholtz-Gl.\ (Frequenzbereich)}{%
      \Delta\underline{\vec H}-(\kappa\mu\, j\omega-\varepsilon\mu\,\omega^2)\underline{\vec H}=0}} &
  \parbox[t]{\linewidth}{%
    \formelblock{Isolator / Vakuum ($\kappa=0$)}{%
      \Delta\vec H=\varepsilon\mu\tfrac{\partial^2\vec H}{\partial t^2}
        \quad(\text{Wellen})}\\[0.25em]%
    \formelblock{Sehr guter Leiter ($\kappa\gg0$)}{%
      \Delta\vec H=\kappa\mu\tfrac{\partial\vec H}{\partial t}
        \quad(\text{Diffusion, Skineff.})}\\[0.15em]%
    \infoblock{$\varrho=0$ im ladungsfreien Raum
      $\Rightarrow$ rechte Seite $=0$.}} \\[0.3em]
\end{formelreihe}
\bereichende

% 6.2  Ebene Welle — Fortpflanzungskonstante
\bereich[cKap6]{6.2 Ebene Welle — Fortpflanzungskonstante}
\begin{formelreihe}
  \parbox[t]{\linewidth}{%
    \infoblock{\textbf{Ebene Welle:} Flächen konst.\ Phase = Ebenen $\perp$
      Ausbreitung ($z$). Es gilt:\\[0.15em]
      \textbullet\ $\vec E\perp\vec H$, beide $\perp$ Ausbreitung\\
      \textbullet\ keine Feldkomp.\ in Ausbreitungsricht.\\
      \textbullet\ bei $\kappa=0$: $\vec E,\vec H$ gleichphasig}\\[0.2em]%
    \formelblock{Ansatz}{%
      \underline{\vec E}=\underline{\vec E}_1 e^{-\underline\gamma z}+\underline{\vec E}_2 e^{+\underline\gamma z}}} &
  \parbox[t]{\linewidth}{%
    \merkblock[cKap6]{Fortpflanzungskonstante}{%
      \underline\gamma=\sqrt{j\omega\mu(\kappa+j\omega\varepsilon)}=\alpha+j\beta}\\[0.25em]%
    \formelblock{Dämpfungskonstante $[\alpha]=\mathrm{Np/m}$}{%
      \alpha=\omega\sqrt{\tfrac{\mu\varepsilon}{2}\Big(\sqrt{1+\tfrac{\kappa^2}{\omega^2\varepsilon^2}}-1\Big)}}\\[0.2em]%
    \formelblock{Phasenkonstante $[\beta]=\mathrm{rad/m}$}{%
      \beta=\omega\sqrt{\tfrac{\mu\varepsilon}{2}\Big(\sqrt{1+\tfrac{\kappa^2}{\omega^2\varepsilon^2}}+1\Big)}}} &
  \parbox[t]{\linewidth}{%
    \merkblock{Verlustloses Dielektrikum ($\kappa=0$)}{%
      \begin{aligned}
        \alpha&=0\\[0.1em]
        \beta&=\omega\sqrt{\mu\varepsilon}=\omega\sqrt{\mu_0\mu_r\varepsilon_0\varepsilon_r}\\[0.1em]
        &=\tfrac{\omega}{c_0}\sqrt{\mu_r\varepsilon_r}
      \end{aligned}}\\[0.2em]%
    \infoblock{$\alpha$: Dämpfung, $\beta$: Phasendrehung.\\
      Verlustlos $\Rightarrow$ keine Dämpfung.}} \\[0.3em]
\end{formelreihe}
\bereichende

% 6.3  Geschwindigkeit, Wellenlänge, Dispersion
\bereich[cKap6]{6.3 Ausbreitungsgeschwindigkeit, Wellenlänge, Dispersion}
\begin{formelreihe}
  \parbox[t]{\linewidth}{%
    \merkblock[cKap6]{Lichtgeschw.\ / Vakuum}{%
      c_0=\tfrac{1}{\sqrt{\mu_0\varepsilon_0}}\approx3\cdot10^8\,\tfrac{\mathrm{m}}{\mathrm{s}}}\\[0.25em]%
    \formelblock{Phasengeschwindigkeit}{%
      v_{ph}=\tfrac{\lambda}{T}=\lambda f=\tfrac{\omega}{\beta}
        =\tfrac{c_0}{\sqrt{\mu_r\varepsilon_r}}}} &
  \parbox[t]{\linewidth}{%
    \merkblock{Wellenlänge}{%
      \lambda=\tfrac{2\pi}{\beta}=\tfrac{c_0}{f}\cdot\tfrac{1}{\sqrt{\mu_r\varepsilon_r}}}\\[0.25em]%
    \infoblock{$\lambda$ ist material\emph{ab}hängig: in Materie
      ($\varepsilon_r,\mu_r\!>\!1$) stets \emph{kleiner} als im Vakuum.}} &
  \parbox[t]{\linewidth}{%
    \formelblock{Gruppengeschwindigkeit}{%
      v_g=\tfrac{1}{\mathrm{d}\beta/\mathrm{d}\omega}
        =\tfrac{v_{ph}}{1-\tfrac{f}{v_{ph}}\tfrac{\mathrm{d}v_{ph}}{\mathrm{d}f}}}\\[0.2em]%
    \infoblock{$v_g\!\neq\!v_{ph}$, falls $v_{ph}=f(\omega)$
      $\Rightarrow$ \emph{Dispersion}. $v_g$ = Geschw.\ der
      Energie-/Informationsübertragung.}} \\[0.3em]
\end{formelreihe}
\bereichende

% 6.4  Feldwellenwiderstand & Poynting
\bereich[cKap6]{6.4 Feldwellenwiderstand \& Energie (Poynting)}
\begin{formelreihe}
  \parbox[t]{\linewidth}{%
    \formelblock{Feldwellenwiderstand (allgemein)}{%
      \underline Z_F=\tfrac{\underline E_{\mathrm{trans}}}{\underline H_{\mathrm{trans}}}
        =\sqrt{\tfrac{j\omega\mu}{\kappa+j\omega\varepsilon}}}\\[0.2em]%
    \formelblock{Dielektrikum ($\kappa=0$, reell)}{%
      Z_F=\sqrt{\tfrac{\mu}{\varepsilon}}=Z_{F0}\sqrt{\tfrac{\mu_r}{\varepsilon_r}}}} &
  \parbox[t]{\linewidth}{%
    \merkblock[cKap6]{Freiraum-Wellenwiderstand}{%
      Z_{F0}=\sqrt{\tfrac{\mu_0}{\varepsilon_0}}=120\pi\,\Omega\approx377\,\Omega}\\[0.2em]%
    \infoblock{$\kappa\!\neq\!0$: $\underline Z_F$ komplex
      $\Rightarrow$ Phasenverschiebung zw.\ $E$ und $H$.}} &
  \parbox[t]{\linewidth}{%
    \merkblock{Poynting-Vektor (Leistungsfluss)}{%
      \vec S=\vec E\times\vec H\;,\quad[S]=\tfrac{\mathrm{W}}{\mathrm{m}^2}}\\[0.2em]%
    \formelblock{Energiedichte}{%
      \mathrm{d}w=\vec E\,\mathrm{d}\vec D+\vec H\,\mathrm{d}\vec B}\\[0.15em]%
    \formelblock{Energieerhaltung}{%
      \tfrac{\partial W}{\partial t}=\underbrace{-\oint\vec S\,\mathrm{d}\vec A}_{\text{Strahlung}}
        \underbrace{-\int\kappa E^2\,\mathrm{d}V}_{\text{Verluste}}}} \\[0.3em]
\end{formelreihe}
\bereichende

% 6.5  Polarisation
\bereich[cKap6]{6.5 Polarisation (bezieht sich auf $\vec E$)}
\begin{formelreihe}
  \parbox[t]{\linewidth}{%
    \formelblock{Lineare Polarisation}{%
      \text{Endpunkt von }\vec E\text{ auf einer Geraden}}\\[0.15em]%
    \infoblock{Bez.\ Erdoberfläche: \emph{horizontal} / \emph{vertikal}.
      $\vec E$ liegt in der Phasenfront-Ebene.}} &
  \parbox[t]{\linewidth}{%
    \formelblock{Elliptische Polarisation}{%
      \text{Endpunkt von }\vec E\text{ beschreibt Ellipse}}\\[0.15em]%
    \infoblock{Ellipse $\to$ Kreis: \emph{zirkulare} Pol.\\
      links-/rechtsdrehend je nach Umlaufsinn
      (in Ausbreitungsrichtung betrachtet).}} &
  \parbox[t]{\linewidth}{%
    \infoblock{Eine ebene Welle setzt sich zusammen aus $E_x,H_y$
      und/oder $E_y,H_x$. Überlagerung zweier senkrechter,
      phasenverschobener Linearwellen $\Rightarrow$ ellipt./zirk.\ Pol.}\\[0.2em]
    \grafik[scale=0.8]{%
      \draw[vec,cKap6] (0,0)--(30:0.9);
      \draw[vec,cKap6] (0,0)--(90:0.9);
      \draw[vec,cKap6] (0,0)--(150:0.9);
      \draw[cKap6,line width=0.3pt] (0.9,0) arc (0:180:0.9);
      \node[font=\tiny] at (0,-0.18){zirkular};}} \\[0.3em]
\end{formelreihe}
\bereichende

% 6.6  Senkrechter Einfall auf Metall
\bereich[cKap6]{6.6 Grenzfläche: senkrechter Einfall auf Metall ($\kappa\to\infty$)}
\begin{formelreihe}
  \parbox[t]{\linewidth}{%
    \formelblock{Randbedingung ideales Metall}{%
      \vec E_t=0\;\Rightarrow\;E_{xh}=-E_{xr}\;,\quad H_{yh}=H_{yr}}\\[0.2em]%
    \infoblock{Hin- + rücklaufende Welle überlagern sich zur
      \emph{stehenden Welle} (vollständige Reflexion).}} &
  \parbox[t]{\linewidth}{%
    \merkblock[cKap6]{Stehende Welle ($z=0$ an Metallfl.)}{%
      \begin{aligned}
        E_x&=2E_{xh}\,\sin\omega t\,\sin\beta z\\[0.1em]
        H_y&=2H_{yh}\,\cos\omega t\,\cos\beta z
      \end{aligned}}} &
  \parbox[t]{\linewidth}{%
    \infoblock{$\vec E$ und $\vec H$ um $90^\circ$ phasenverschoben
      (zeitlich \emph{und} örtlich) $\Rightarrow$
      \textbf{kein} Leistungsfluss (Leistungsbilanz $=0$).\\[0.15em]
      Knoten von $E$ und $H$ um $\lambda/4$ versetzt.\\[0.15em]
      Anwendung: Parabolantenne (\glqq Sat-Schüssel\grqq).}} \\[0.3em]
\end{formelreihe}
\bereichende

% 6.7  Senkrechter Einfall auf Dielektrikum
\bereich[cKap6]{6.7 Grenzfläche: senkrechter Einfall auf Dielektrikum ($\mu_r=1$)}
\begin{formelreihe}
  \parbox[t]{\linewidth}{%
    \merkblock[cKap6]{Reflexionsfaktor (Feldstärke)}{%
      r_e=r_m=\dfrac{\sqrt{\varepsilon_{r2}}-\sqrt{\varepsilon_{r1}}}{\sqrt{\varepsilon_{r2}}+\sqrt{\varepsilon_{r1}}}}\\[0.25em]%
    \formelblock{Transmissionsfaktoren}{%
      t_e=\tfrac{2}{1+\sqrt{\varepsilon_{r2}/\varepsilon_{r1}}}\;,\quad
      t_m=\tfrac{2}{1+\sqrt{\varepsilon_{r1}/\varepsilon_{r2}}}}} &
  \parbox[t]{\linewidth}{%
    \formelblock{Leistungsbezogen}{%
      \begin{aligned}
        r&=-\sqrt{r_e r_m}=\tfrac{\sqrt{\varepsilon_{r1}}-\sqrt{\varepsilon_{r2}}}{\sqrt{\varepsilon_{r2}}+\sqrt{\varepsilon_{r1}}}\\[0.15em]
        t&=\sqrt{t_e t_m}=\tfrac{2\sqrt[4]{\varepsilon_{r1}\varepsilon_{r2}}}{\sqrt{\varepsilon_{r2}}+\sqrt{\varepsilon_{r1}}}
      \end{aligned}}\\[0.15em]%
    \merkblock{Leistungsbilanz}{r_e r_m+t_e t_m=1}} &
  \parbox[t]{\linewidth}{%
    \infoblock{$0<t_e\le2$,\; $-1<r_e<1$.\\
      $r$ negativ $\Rightarrow$ $180^\circ$-Phasensprung des $\vec E$
      (bei Übergang ins dichtere Medium $\varepsilon_{r2}>\varepsilon_{r1}$).\\[0.2em]
      Reflekt.\ = durchgel.\ Leistung erst bei
      $\varepsilon_{r2}/\varepsilon_{r1}\approx34$.\\[0.15em]
      Glas ($\varepsilon_r\!\approx\!2{,}2$): $\approx95\%$ Transmission.}} \\[0.3em]
\end{formelreihe}
\bereichende

% 6.8  Schräger Einfall
\bereich[cKap6]{6.8 Schräger Einfall — Brechung, Totalreflexion, Brewster}
\begin{formelreihe}
  \parbox[t]{\linewidth}{%
    \merkblock[cKap6]{Brechungsgesetz ($n=\sqrt{\varepsilon_r}$)}{%
      \tfrac{\sin\beta}{\sin\alpha}=\sqrt{\tfrac{\varepsilon_{r1}}{\varepsilon_{r2}}}=\tfrac{n_1}{n_2}}\\[0.2em]%
    \infoblock{$\alpha$ = Einfalls-/Reflexionswinkel,
      $\beta$ = Brechungswinkel. Welle wird im dichteren
      Medium \glqq abgebremst\grqq\ ($v_{ph}=c_0/\sqrt{\varepsilon_r}$).}} &
  \parbox[t]{\linewidth}{%
    \merkblock{Totalreflexion (nur $\varepsilon_{r1}>\varepsilon_{r2}$)}{%
      \alpha_g=\arcsin\sqrt{\tfrac{\varepsilon_{r2}}{\varepsilon_{r1}}}}\\[0.2em]%
    \infoblock{Für $\alpha>\alpha_g$ (vom dichten ins dünne Medium):
      Welle wird komplett reflektiert.\\
      Anwendung: \emph{Glasfaser}.}} &
  \parbox[t]{\linewidth}{%
    \merkblock[cKap6]{Brewster-Winkel}{%
      \alpha_B=\arctan\sqrt{\tfrac{\varepsilon_{r2}}{\varepsilon_{r1}}}}\\[0.2em]%
    \infoblock{Bei Einfall unter $\alpha_B$ wird der reflektierte
      Anteil $=0$ $\Rightarrow$ Welle tritt vollständig über.\\
      Anwendung: \emph{Brewster-Fenster} (Gas-Laser).}} \\[0.3em]
\end{formelreihe}
\bereichende

%% ---- EDy7: Wellenleiter ----------------------------------

% 7.1  Koaxialkabel (TEM)
\bereich[cKap7]{7.1 Koaxialkabel (TEM-Welle)}
\begin{formelreihe}
  \parbox[t]{\linewidth}{%
    \formelblock{Felder ($d$ innen, $D$ außen)}{%
      \vec E=\tfrac{U}{r\,\ln(D/d)}\vec e_r\;,\quad
      \vec H=\tfrac{I}{2\pi r}\vec e_\varphi}\\[0.2em]%
    \merkblock[cKap7]{Leitungswellenwiderstand}{%
      Z_L=Z_{F0}\tfrac{\ln(D/d)}{2\pi\sqrt{\varepsilon_r}}
        =\tfrac{60\,\Omega}{\sqrt{\varepsilon_r}}\ln\tfrac{D}{d}}} &
  \parbox[t]{\linewidth}{%
    \formelblock{Dämpfung — Leiterverluste}{%
      \alpha_L=\sqrt{\tfrac{f\mu\varepsilon_r}{\pi\kappa}}
        \cdot\tfrac{1}{120\,\Omega}\cdot\tfrac{1}{D}\cdot\tfrac{1+D/d}{\ln(D/d)}}\\[0.2em]%
    \formelblock{Dämpfung — dielektr.\ ($\tan\delta$ klein)}{%
      \alpha_D\approx\tfrac{\pi f}{c_0}\sqrt{\varepsilon_r}\,\tan\delta}\\[0.15em]%
    \infoblock{Niedrige $f$: Leiterverluste dominieren;
      ab GHz: dielektr.\ Verluste $\to$ Luft/Schaum.}} &
  \parbox[t]{\linewidth}{%
    \merkblock{Optimum (min.\ Dämpfung)}{%
      (D/d)_{opt}=3{,}59\;\Rightarrow\;Z_L=\tfrac{77\,\Omega}{\sqrt{\varepsilon_r}}}\\[0.2em]%
    \infoblock{$\to$ $Z_L\!\approx\!50\,\Omega$ (PE, $\varepsilon_r\!=\!2{,}2$),
      $75\,\Omega$ (Luft).\\
      $60\,\Omega/\!\sqrt{\varepsilon_r}$: max.\ Spannungsfestigkeit,
      $30\,\Omega/\!\sqrt{\varepsilon_r}$: max.\ Leistung.}\\[0.15em]%
    \formelblock{Obere Grenzfrequenz (Eindeutigkeit)}{%
      f_g\approx\tfrac{0{,}174\,\text{GHz}\cdot\text{m}}{D\sqrt{\varepsilon_r}}}} \\[0.3em]
\end{formelreihe}
\bereichende

% 7.2  Hohlleiter — Grundlagen
\bereich[cKap7]{7.2 Hohlleiter — Grundlagen (Hochpass-Charakter)}
\begin{formelreihe}
  \parbox[t]{\linewidth}{%
    \infoblock{Hohles Metallrohr. Kein TEM möglich (nur mit Innenleiter).
      Es existieren:\\[0.15em]
      \textbullet\ \textbf{H-Welle} (TE): $H_z\neq0,\;E_z=0$\\
      \textbullet\ \textbf{E-Welle} (TM): $E_z\neq0,\;H_z=0$}\\[0.2em]%
    \merkblock[cKap7]{Grenzfrequenz / cut-off}{%
      f_c=\tfrac{c}{\lambda_c}\;,\quad c=\tfrac{c_0}{\sqrt{\varepsilon_r\mu_r}}}} &
  \parbox[t]{\linewidth}{%
    \formelblock{Durchlass (Ausbreitung)}{%
      f>f_c\;\Leftrightarrow\;\lambda<\lambda_c\;:\;
      \beta=\tfrac{2\pi}{\lambda}\sqrt{1-\big(\tfrac{\lambda}{\lambda_c}\big)^2}}\\[0.2em]%
    \formelblock{Sperrbereich (Dämpfung)}{%
      f<f_c\;:\;\alpha=\tfrac{2\pi}{\lambda}\sqrt{\big(\tfrac{\lambda}{\lambda_c}\big)^2-1}}\\[0.15em]%
    \merkblock{Hohlleiterwellenlänge}{%
      \lambda_H=\tfrac{\lambda}{\sqrt{1-(\lambda/\lambda_c)^2}}}} &
  \parbox[t]{\linewidth}{%
    \formelblock{Phasen-/Gruppengeschwindigkeit}{%
      \begin{aligned}
        v_{ph}&=\tfrac{f\lambda}{\sqrt{1-(\lambda/\lambda_c)^2}}>c\\[0.1em]
        v_g&=f\lambda\sqrt{1-(\lambda/\lambda_c)^2}<c
      \end{aligned}}\\[0.15em]%
    \infoblock{$v_{ph}\cdot v_g=c^2$. Anwendung Dämpfung:
      \emph{Wabenkamin} (Lüftung geschirmter Räume).}} \\[0.3em]
\end{formelreihe}
\bereichende

% 7.3  Rechteckhohlleiter — Hmn/Emn
\bereich[cKap7]{7.3 Rechteckhohlleiter — H$_{mn}$/E$_{mn}$-Wellen}
\begin{formelreihe}
  \parbox[t]{\linewidth}{%
    \merkblock[cKap7]{Grenzwellenlänge (Breite $a$, Höhe $b$)}{%
      \lambda_c=\tfrac{2}{\sqrt{(m/a)^2+(n/b)^2}}}\\[0.2em]%
    \formelblock{Phasenkonstante}{%
      \beta=\sqrt{\omega^2\mu\varepsilon-(\tfrac{m\pi}{a})^2-(\tfrac{n\pi}{b})^2}}\\[0.15em]%
    \infoblock{$m$: Max./Nullst.\ über $a$, $n$: über $b$.
      $f_c$ hängt nur von $m,n$ ab (E \& H mit gleichen Indizes
      $\Rightarrow$ gleiches $f_c$).}} &
  \parbox[t]{\linewidth}{%
    \formelblock{Grundwelle H$_{10}$}{%
      \lambda_c=2a\;,\qquad f_c=\tfrac{c}{2a}}\\[0.2em]%
    \infoblock{H-Wellen: auch H$_{01}$, H$_{20}$, H$_{11}$ \dots\\
      E-Wellen: \textbf{kein} E$_{m0}$/E$_{0n}$
      (dann $E_z\!=\!0$) $\Rightarrow$ Grundwelle \textbf{E$_{11}$}.}\\[0.15em]%
    \formelblock{Feldwellenwiderstände}{%
      \begin{aligned}
        Z_{FH}&=\tfrac{Z_{F0}}{\sqrt{1-(\lambda/\lambda_c)^2}}\\[0.1em]
        Z_{FE}&=Z_{F0}\sqrt{1-(\lambda/\lambda_c)^2}
      \end{aligned}}} &
  \parbox[t]{\linewidth}{%
    \merkblock{Eindeutigkeitsbereich ($a\ge2b$)}{%
      f_{cH10}<f<f_{cH20}=2f_{cH10}}\\[0.2em]%
    \infoblock{Nur \emph{ein} Modus (H$_{10}$) ausbreitungsfähig
      $\Rightarrow$ keine Modenverzerrung.\\
      Praktisch genutzt (Sicherheitsabstand):}\\[0.15em]%
    \merkblock[cKap7]{Betriebsfrequenzbereich}{%
      1{,}25\,f_{cH10}<f<1{,}9\,f_{cH10}}} \\[0.3em]
\end{formelreihe}
\bereichende

% 7.4  H10-Welle + Leistung (klausurrelevant)
\bereich[cKap7]{7.4 H$_{10}$-Grundwelle — Feldbild \& übertragene Leistung \; (klausurrelevant!)}
\begin{formelreihe}
  \parbox[t]{\linewidth}{%
    \formelblock{Feldkomponenten (nur $E_y,H_x,H_z$)}{%
      \begin{aligned}
        E_y&=E_0\,\sin\!\tfrac{\pi x}{a}\,e^{-j\beta z}\\[0.1em]
        H_x&=-\tfrac{E_y}{Z_{FH}}\;,\quad H_z\sim\cos\tfrac{\pi x}{a}
      \end{aligned}}\\[0.15em]%
    \grafik[scale=0.9]{%
      \draw[line width=0.8pt] (0,0) rectangle (2.0,1.0);
      \foreach \x in {0.2,0.45,...,1.85}{%
        \pgfmathsetmacro\h{0.86*sin(90*\x)}
        \draw[vec,blue!70!black] (\x,0.07)--(\x,{0.07+\h});}
      \draw[<->,line width=0.3pt] (0,-0.12)--(2.0,-0.12) node[midway,below,font=\tiny]{$a$};
      \draw[<->,line width=0.3pt] (-0.12,0)--(-0.12,1.0) node[midway,left,font=\tiny]{$b$};
      \node[font=\tiny,blue!70!black] at (1.0,1.16){$\vec E=E_y$};}} &
  \parbox[t]{\linewidth}{%
    \formelblock{Wirkleistungsdichte in $z$}{%
      \langle S_z\rangle=\tfrac{E_0^2}{2Z_{FH}}\sin^2\!\tfrac{\pi x}{a}}\\[0.2em]%
    \merkblock[cKap7]{Übertragene Leistung H$_{10}$}{%
      P=\dfrac{a\,b\,E_0^2}{4\,Z_{FH}}
        =\dfrac{a\,b\,E_0^2}{4\,Z_{F0}}\sqrt{1-\big(\tfrac{\lambda}{2a}\big)^2}}\\[0.15em]%
    \infoblock{$E_0$ = Amplitude in Rohrmitte.
      $P\!\to\!0$ bei $f\!\to\!f_c$ ($Z_{FH}\!\to\!\infty$).}} &
  \parbox[t]{\linewidth}{%
    \merkblock{Maximale (Durchbruch-)Leistung}{%
      P_{max}=\tfrac{a\,b\,E_D^2}{4\,Z_{F0}}\sqrt{1-\big(\tfrac{\lambda}{2a}\big)^2}}\\[0.2em]%
    \formelblock{Wanddämpfung H$_{10}$}{%
      \alpha_{Wand}=\tfrac{2R_F}{Z_{F0}}\cdot
        \tfrac{a/2b+(\lambda_0/2a)^2}{a\sqrt{1-(\lambda_0/2a)^2}}}\\[0.15em]%
    \infoblock{$E_D$ = Durchbruchfeldstärke.
      $R_F=\tfrac{1}{\kappa\delta}$ (Skineffekt) steigt mit $f$.}} \\[0.3em]
\end{formelreihe}
\bereichende

% 7.5  Rundhohlleiter
\bereich[cKap7]{7.5 Rundhohlleiter (Zylinderkoord., Besselfunktionen)}
\begin{formelreihe}
  \parbox[t]{\linewidth}{%
    \merkblock[cKap7]{Grenzwellenlängen (Durchm.\ $D$)}{%
      \begin{aligned}
        \text{H}_{mn}:\;&\lambda_{c}=\tfrac{\pi D}{\eta'_{mn}}\\[0.1em]
        \text{E}_{mn}:\;&\lambda_{c}=\tfrac{\pi D}{\eta_{mn}}
      \end{aligned}}\\[0.15em]%
    \infoblock{$\eta'_{mn}$: $n$-te Nullst.\ der \emph{Ableitung} der
      Besselfkt.\ $m$-ter Ordnung; $\eta_{mn}$: der Besselfkt.\ selbst.}} &
  \parbox[t]{\linewidth}{%
    \formelblock{Nullstellen (klein $\to$ groß)}{%
      \begin{array}{@{}l@{\;}l@{}}
        \text{H}: & \eta'_{11}\!=\!1{,}84,\ \eta'_{21}\!=\!3{,}05,\ \eta'_{01}\!=\!3{,}83\\
        \text{E}: & \eta_{01}\!=\!2{,}40,\ \eta_{11}\!=\!3{,}83,\ \eta_{21}\!=\!5{,}13
      \end{array}}\\[0.2em]%
    \merkblock{Grundwelle}{%
      \text{H}_{11}\;(\eta'_{11}=1{,}84,\ \text{kleinstes }f_c)}} &
  \parbox[t]{\linewidth}{%
    \formelblock{Eindeutigkeitsbereich}{%
      1{,}306\,D<\lambda<1{,}706\,D}\\[0.2em]%
    \infoblock{$m$: Max.\ über halben Umfang, $n$: Nullst.\ über Radius.\\[0.15em]
      H$_{01}$ hat \emph{geringere} Dämpfung als H$_{11}$
      (nimmt mit $f$ ab!) $\Rightarrow$ verlustarme Fernübertragung.
      Rundhohlleiter: kleinere Wandverluste als Rechteck.}} \\[0.3em]
\end{formelreihe}
\bereichende

% 7.6  Hohlraumresonatoren
\bereich[cKap7]{7.6 Hohlraumresonatoren}
\begin{formelreihe}
  \parbox[t]{\linewidth}{%
    \infoblock{Beidseitig leitend abgeschlossener Hohlleiter
      ($d=p\cdot\lambda_H/2$) wirkt wie $LC$-Schwingkreis.}\\[0.2em]%
    \merkblock[cKap7]{Resonanzfrequenzen (Quader $a\!\times\!b\!\times\!d$)}{%
      f_{mnp}=\tfrac{c_0}{2\sqrt{\varepsilon_r\mu_r}}
        \sqrt{(\tfrac{m}{a})^2+(\tfrac{n}{b})^2+(\tfrac{p}{d})^2}}} &
  \parbox[t]{\linewidth}{%
    \infoblock{Güten bis $\approx10000$, nutzbar bis $100$\,GHz.
      Anregung nur über Koppelöffnungen (Schlitze/Löcher).}\\[0.2em]%
    \formelblock{Anwendungen}{%
      \begin{aligned}
        &\text{Duplexfilter (GSM/UMTS-Basisstation)}\\
        &\text{Frequenzmesser (variabler Hohlraum)}
      \end{aligned}}\\[0.15em]%
    \infoblock{Vorteil ggü.\ $LC$: hohe Güte im GHz-Bereich
      ($LC$ dort nicht mehr baubar). Nachteil: größer/schwerer.}} &
  \parbox[t]{\linewidth}{%
    \merkblock{EMV-Merksatz}{%
      \text{Jeder geschl.\ leitf.\ Hohlraum hat Resonanzfrequenzen}}\\[0.2em]%
    \infoblock{PC-Gehäuse: interne Taktfreq.\ regen Gehäuseresonanzen
      an $\Rightarrow$ Feldüberhöhung ($\sim$Güte) $\Rightarrow$ Abstrahlung.\\[0.15em]
      Abhilfe: viele \emph{kleine} Löcher statt weniger großer,
      Kontaktfedern, geschirmte Kabel, Ferritkerne/Netzfilter.}} \\[0.3em]
\end{formelreihe}
\bereichende

% 7.7  Leitungstheorie — Telegrafengleichungen
\bereich[cKap7]{7.7 Leitungstheorie — Telegrafengleichungen \; (klausurrelevant!)}
\begin{formelreihe}
  \parbox[t]{\linewidth}{%
    \formelblock{Leitungsbeläge (längenbezogen)}{%
      R',L'\ (\text{längs}),\quad G',C'\ (\text{quer})}\\[0.2em]%
    \merkblock[cKap7]{Leitungswellenwiderstand}{%
      \underline Z_L=\sqrt{\tfrac{R'+j\omega L'}{G'+j\omega C'}}}\\[0.15em]%
    \formelblock{Fortpflanzungskonstante}{%
      \underline\gamma=\alpha+j\beta=\sqrt{(R'+j\omega L')(G'+j\omega C')}}} &
  \parbox[t]{\linewidth}{%
    \merkblock{Verlustlose Leitung ($R'=G'=0$)}{%
      \begin{aligned}
        Z_L&=\sqrt{\tfrac{L'}{C'}}\quad(\text{reell})\\[0.1em]
        \beta&=\omega\sqrt{L'C'}\\[0.1em]
        v_{ph}&=\tfrac{1}{\sqrt{L'C'}}=\tfrac{c_0}{\sqrt{\varepsilon_r\mu_r}}
      \end{aligned}}} &
  \parbox[t]{\linewidth}{%
    \formelblock{Telegrafengl.\ ($z$ vom Ende gezählt)}{%
      \begin{aligned}
        \underline U(z)&=\underline U(0)\cos\beta z+jZ_L\underline I(0)\sin\beta z\\[0.1em]
        \underline I(z)&=\underline I(0)\cos\beta z+j\tfrac{\underline U(0)}{Z_L}\sin\beta z
      \end{aligned}}\\[0.15em]%
    \infoblock{Kurzschluss ($U(0)\!=\!0$) / Leerlauf
      ($I(0)\!=\!0$): stehende Wellen, $U,I$ um $\lambda/4$ versetzt.}} \\[0.3em]
\end{formelreihe}
\bereichende

% 7.8  Leitungstransformation & Reflexion
\bereich[cKap7]{7.8 Leitungstransformation \& Reflexionsfaktor \; (klausurrelevant!)}
\begin{formelreihe}
  \parbox[t]{\linewidth}{%
    \formelblock{Impedanztransformation (verlustlos)}{%
      \underline Z(z)=Z_L\cdot
        \tfrac{\underline Z(0)+jZ_L\tan\beta z}{Z_L+j\underline Z(0)\tan\beta z}}\\[0.2em]%
    \merkblock[cKap7]{$\lambda/4$-Transformator}{%
      \underline Z(\tfrac{\lambda}{4})=\tfrac{Z_L^2}{\underline Z(0)}}\\[0.15em]%
    \infoblock{$z=n\tfrac{\lambda}{2}$: $Z(z)=Z(0)$ (keine Transf.).}} &
  \parbox[t]{\linewidth}{%
    \formelblock{$\lambda/4$-Sonderfälle}{%
      \begin{aligned}
        \text{Leerlauf}&\to\text{Kurzschluss}\\
        \text{Kurzschluss}&\to\text{Leerlauf}\\
        C&\to L=C Z_L^2\\
        L&\to C=L/Z_L^2
      \end{aligned}}\\[0.15em]%
    \infoblock{$\Rightarrow$ Blindelemente aus Leitungsstücken.}} &
  \parbox[t]{\linewidth}{%
    \merkblock{Reflexionsfaktor}{%
      \underline r(z)=\tfrac{\underline Z(0)-Z_L}{\underline Z(0)+Z_L}\,e^{-2j\beta z}}\\[0.2em]%
    \infoblock{$|r|$ ändert sich längs (verlustloser) Leitung nicht.\\
      \textbf{Anpassung} $\underline Z(0)=Z_L$: $r=0$, keine Reflexion,
      keine Transformation.\\[0.15em]
      Rückflussdämpfung $=-20\lg|r|$ [dB].\\
      $\Rightarrow$ lange Digitalleitungen mit $Z_L$ abschließen!}} \\[0.3em]
\end{formelreihe}
\bereichende

%% ---- EDy8: Antennen --------------------------------------

% 8.1  Ersatzschaltbild & Wirkungsgrad
\bereich[cKap8]{8.1 Antenne — Ersatzschaltbild \& Wirkungsgrad}
\begin{formelreihe}
  \parbox[t]{\linewidth}{%
    \infoblock{Antenne wandelt Leitungswelle $\leftrightarrow$ Freiraumwelle
      ($Z_L\!\leftrightarrow\!Z_{F0}$). Passive Antennen sind
      \emph{reziprok} (Sende- = Empfangsverhalten).}\\[0.2em]%
    \formelblock{Antennenimpedanz}{%
      \underline Z_A=\underbrace{R_V}_{\text{Verlust}}
        +\underbrace{R_S}_{\text{Strahlung}}+jX_A}} &
  \parbox[t]{\linewidth}{%
    \formelblock{Abgestrahlte / Verlustleistung}{%
      P_S=\tfrac{1}{2}\hat\imath_A^2 R_S\;,\quad
      P_V=\tfrac{1}{2}\hat\imath_A^2 R_V}\\[0.25em]%
    \merkblock[cKap8]{Wirkungsgrad}{%
      \eta=\tfrac{P_S}{P_S+P_V}=\tfrac{R_S}{R_S+R_V}}} &
  \parbox[t]{\linewidth}{%
    \infoblock{$R_V$ hängt ab von: Oberflächenwiderstand
      (Skineffekt), dielektr.\ Verluste, Erdströme.\\[0.2em]
      Bauformen: Dipol, Stab ($\lambda/4$), Yagi, Helix,
      Loop/Ferrit, Horn, Parabol, Patch, Schlitz.}} \\[0.3em]
\end{formelreihe}
\bereichende

% 8.2  Elementardipol
\bereich[cKap8]{8.2 Elementardipol (Hertzscher Dipol, $l\ll\lambda$)}
\begin{formelreihe}
  \parbox[t]{\linewidth}{%
    \merkblock[cKap8]{Fernfeld ($r\gg\lambda$)}{%
      \begin{aligned}
        E_\vartheta&=\tfrac{\hat\imath\,l}{2\varepsilon_0 c_0\lambda}
          \cdot\tfrac{\sin\vartheta}{r}\;,\quad E_r=0\\[0.1em]
        H_\varphi&=\tfrac{\hat\imath\,l}{2\lambda}\cdot\tfrac{\sin\vartheta}{r}
      \end{aligned}}\\[0.15em]%
    \infoblock{$[E]=$ V/m, $[H]=$ A/m. $E,H$ in Phase $\Rightarrow$ Kugelwelle,
      Wirkleistung radial nach außen.}} &
  \parbox[t]{\linewidth}{%
    \formelblock{Feldwellenwiderstand (Fernfeld)}{%
      Z_F=\tfrac{E_\vartheta}{H_\varphi}=\tfrac{1}{\varepsilon_0 c_0}
        =Z_{F0}=120\pi\,\Omega\approx377\,\Omega}\\[0.25em]%
    \formelblock{Nahfeld ($r\ll\lambda$)}{%
      \underline Z_{F,nah}=-j\tfrac{\lambda}{2\pi r}Z_{F0}}\\[0.15em]%
    \infoblock{Nahfeld: $E,H$ um $90^\circ$ verschoben
      $\Rightarrow$ Blindleistungsfeld.}} &
  \parbox[t]{\linewidth}{%
    \merkblock{Nah-/Fernfeld-Grenze}{%
      r<\tfrac{\lambda}{4\pi}:\text{Nahfeld}\;,\quad
      r>\tfrac{\lambda}{\pi}:\text{Fernfeld}}\\[0.2em]%
    \infoblock{\emph{Elektrischer} Strahler (Dipol): im Nahfeld
      \emph{kapazitiv}.\\
      \emph{Magnetischer} Strahler (Schleife, $\pi d\!\ll\!\lambda$):
      im Nahfeld \emph{induktiv}.}} \\[0.3em]
\end{formelreihe}
\bereichende

% 8.3  Richtcharakteristik & Kenngrößen
\bereich[cKap8]{8.3 Richtcharakteristik \& Kenngrößen}
\begin{formelreihe}
  \parbox[t]{\linewidth}{%
    \merkblock[cKap8]{Richtcharakteristik (feldnormiert)}{%
      C(\vartheta,\varphi)=\tfrac{E(\vartheta,\varphi)}{E_{max}}
        =\tfrac{H(\vartheta,\varphi)}{H_{max}}=\tfrac{U(\vartheta,\varphi)}{U_{max}}}\\[0.2em]%
    \formelblock{Leistungsbezogen}{%
      D(\vartheta,\varphi)=\tfrac{S_r(\vartheta,\varphi)}{S_{r,max}}=C^2}} &
  \parbox[t]{\linewidth}{%
    \infoblock{\textbf{Begriffe:}\\
      \textbullet\ Hauptstrahlrichtung: Richtung von $E_{max}$\\
      \textbullet\ \emph{Öffnungswinkel} (Halbwertsbr., 3dB-Breite):
      $E$ um $\le3$\,dB abgesunken\\
      \textbullet\ Nebenkeulendämpfung [dB]\\
      \textbullet\ Vor-Rück-Verhältnis / Rückdämpfung}} &
  \parbox[t]{\linewidth}{%
    \infoblock{Diagramme nach Schnittebene / Feldkomponente:\\
      \textbullet\ \emph{E-Ebenen}-Diagramm\\
      \textbullet\ \emph{H-Ebenen}-Diagramm\\[0.2em]
      Wegen hoher Dynamik oft in dB. In dB verschwindet
      Unterschied zwischen $C$ und $D$.}} \\[0.3em]
\end{formelreihe}
\bereichende

% 8.4  Antennengewinn
\bereich[cKap8]{8.4 Antennengewinn}
\begin{formelreihe}
  \parbox[t]{\linewidth}{%
    \merkblock[cKap8]{Gewinn (Hauptstrahlrichtung)}{%
      G=\tfrac{S_{r,max}}{S_{r,max}(\text{Bezugsant.})}\;,\quad
      g=10\,\mathrm{dB}\cdot\lg G}\\[0.2em]%
    \formelblock{Isotroper Kugelstrahler}{%
      S_r=\tfrac{P}{4\pi r^2}\;\big[\tfrac{\mathrm W}{\mathrm m^2}\big]
      \quad(\text{fiktive Bezugsant.})}} &
  \parbox[t]{\linewidth}{%
    \infoblock{Bezugsantennen:\\
      \textbullet\ \textbf{dBi}: isotroper Kugelstrahler\\
      \textbullet\ \textbf{dBd}: Halbwellendipol\\[0.2em]
      Umrechnung: Dipol strahlt $1{,}64\times$ ($+2{,}15$\,dB)
      stärker als Kugelstrahler:\\
      $g_{[\mathrm{dBi}]}=g_{[\mathrm{dBd}]}+2{,}15$\,dB}} &
  \parbox[t]{\linewidth}{%
    \formelblock{Elementardipol}{%
      G=1{,}5\;\mathrel{\widehat{=}}\;1{,}76\,\mathrm{dBi}}\\[0.2em]%
    \infoblock{Energieerhaltung: mehr Gewinn $\Rightarrow$ kleinerer
      Öffnungswinkel (SAT-Antennen: wenige Grad).\\[0.15em]
      ERP/EIRP: fiktive Sendeleistung für gleiche Leistungsdichte
      (Bezug Dipol bzw.\ Kugelstrahler).}} \\[0.3em]
\end{formelreihe}
\bereichende

% 8.5  Wirkfläche, effektive Länge, Strahlungswiderstand
\bereich[cKap8]{8.5 Wirkfläche, effektive Länge, Strahlungswiderstand}
\begin{formelreihe}
  \parbox[t]{\linewidth}{%
    \formelblock{Antennenwirkfläche (Empfang)}{%
      A_w=\tfrac{P_e}{S_r}}\\[0.2em]%
    \merkblock[cKap8]{Zusammenhang mit Gewinn}{%
      A_w=\tfrac{\lambda^2}{4\pi}\,G\;\;[\mathrm m^2]}\\[0.15em]%
    \infoblock{Flächenstrahler (Horn, Parabol): $A_w\!\approx\!A_{geom}$
      $\Rightarrow$ $G=\tfrac{4\pi}{\lambda^2}A$.}} &
  \parbox[t]{\linewidth}{%
    \formelblock{Effektive Länge (Empfang)}{%
      U_0=l_{e\!f\!f}\cdot E}\\[0.2em]%
    \formelblock{aus Wirkfläche}{%
      l_{e\!f\!f}=\sqrt{\tfrac{4R_S\,A_w}{Z_{F0}}}}\\[0.15em]%
    \infoblock{$l_{e\!f\!f}$: reale Stromverteilung umgerechnet
      auf konstante über $l_{e\!f\!f}$.}} &
  \parbox[t]{\linewidth}{%
    \formelblock{Strahlungswiderstand (Def.)}{%
      R_S=\tfrac{2P_S}{\hat\imath^2}}\\[0.2em]%
    \merkblock{Elementardipol}{%
      R_S=80\pi^2\,\Omega\cdot\big(\tfrac{l}{\lambda}\big)^2}\\[0.15em]%
    \infoblock{$l/\lambda\!=\!1/100$: $R_S\!=\!0{,}08\,\Omega$
      $\Rightarrow$ schlechter Wirkungsgrad.}} \\[0.3em]
\end{formelreihe}
\bereichende

% 8.6  Dipol
\bereich[cKap8]{8.6 Dipol ($\lambda/2$, Ganzwelle, $\lambda/4$)}
\begin{formelreihe}
  \parbox[t]{\linewidth}{%
    \merkblock[cKap8]{$\lambda/2$-Dipol (ideal dünn)}{%
      \underline Z_A=R_S+jX_A=73{,}13\,\Omega+j\,42{,}54\,\Omega}\\[0.2em]%
    \infoblock{Induktiv $\Rightarrow$ Kürzen um Faktor $0{,}96$
      $\to$ reell $73\,\Omega$. Reale Dipole: $R_S\!\approx\!50\text{--}60\,\Omega$.}\\[0.15em]%
    \formelblock{Tatsächliche Länge}{%
      l=0{,}96\cdot\tfrac{\lambda}{2}-d}} &
  \parbox[t]{\linewidth}{%
    \formelblock{Mittlerer Wellenwiderstand}{%
      Z_M=120\,\Omega\cdot\ln\!\big(0{,}575\tfrac{l}{d}\big)}\\[0.2em]%
    \formelblock{Güte (Serienschwingkreis)}{%
      Q=\tfrac{1}{R_S}\sqrt{\tfrac{L}{C}}}\\[0.15em]%
    \infoblock{Dickere Dipole: kleinere Güte $\Rightarrow$
      \emph{höhere} Bandbreite (Breitbanddipol).\\[0.15em]
      $\lambda/4$-Strahler über leitender Ebene: $+3$\,dB Gewinn
      (real $\approx-3$\,dBd bei $\kappa<\infty$).}} &
  \parbox[t]{\linewidth}{%
    \formelblock{Richtcharakteristiken $C(\vartheta)$}{%
      \begin{array}{@{}l@{}}
        \text{Elem.dipol: }|\sin\vartheta|\;(90^\circ,1{,}76\,\text{dBi})\\[0.2em]
        \tfrac{\lambda}{2}: \big|\tfrac{\cos(\frac{\pi}{2}\cos\vartheta)}{\sin\vartheta}\big|\;(78^\circ,2{,}15\,\text{dBi})\\[0.2em]
        \text{Ganzw.: }\big|\tfrac{\cos(\pi\cos\vartheta)+1}{2\sin\vartheta}\big|\;(47^\circ,3{,}8\,\text{dBi})
      \end{array}}\\[0.15em]%
    \grafik[scale=0.8]{%
      \draw[cKap8,line width=0.5pt] (0,0.45) circle (0.45);
      \draw[cKap8,line width=0.5pt] (0,-0.45) circle (0.45);
      \draw[dashed,line width=0.2pt] (-0.7,0)--(0.7,0);
      \node[font=\tiny] at (0,-1.05){$|\sin\vartheta|$};}} \\[0.3em]
\end{formelreihe}
\bereichende

% 8.7  Gruppenantennen
\bereich[cKap8]{8.7 Gruppenantennen (Fernfeld)}
\begin{formelreihe}
  \parbox[t]{\linewidth}{%
    \merkblock[cKap8]{Multiplikatives Gesetz}{%
      C_{ges}=C_{E}\cdot C_{Gr}}\\[0.2em]%
    \infoblock{Gesamt-Richtcharakteristik = Einzelstrahler $\times$
      Gruppen-Charakteristik. Voraussetzung: baugleiche Strahler,
      gleicher Abstand $b$, gleiche Amplitude, konst.\
      Phasendifferenz $\delta$.}} &
  \parbox[t]{\linewidth}{%
    \formelblock{Gruppengewinn ($N$ Strahler)}{%
      \sqrt{G_{Gr}}=\tfrac{1}{\sqrt N}
        \left|\tfrac{\sin(N u/2)}{\sin(u/2)}\right|}\\[0.2em]%
    \merkblock{mit}{%
      u=\delta+\tfrac{2\pi}{\lambda_0}\,b\cos\vartheta}\\[0.15em]%
    \infoblock{$\Delta l=n\lambda$: konstruktiv (doppelte Feldst.);
      $\Delta l=(2n\!+\!1)\tfrac{\lambda}{2}$: Auslöschung.}} &
  \parbox[t]{\linewidth}{%
    \infoblock{$\delta=0$ (gleichphasig): Maximum \emph{quer} zur Zeile.\\[0.15em]
      Variable Phase $\delta$ $\Rightarrow$ Hauptkeule schwenkt
      \emph{ohne} mechanische Bewegung (\emph{Phased Array}).\\[0.2em]
      Verdopplung der Antennenzahl: max.\ $+3$\,dB Gewinn.\\[0.15em]
      \emph{Yagi-Uda}: Strahlungskopplung (Reflektor $l\!>\!\tfrac{\lambda}{2}$,
      Direktoren $l\!<\!\tfrac{\lambda}{2}$).}} \\[0.3em]
\end{formelreihe}
\bereichende

% 8.8  Unerwünschte Abstrahlung (EMV)
\bereich[cKap8]{8.8 Unerwünschte Abstrahlung (EMV)}
\begin{formelreihe}
  \parbox[t]{\linewidth}{%
    \merkblock[cKap8]{Empfindlichkeit Funkempfänger}{%
      P_E\approx1\,\mathrm{pW}\;\mathrel{\widehat{=}}\;-90\,\mathrm{dBm}}\\[0.2em]%
    \formelblock{Nötige Leistungsdichte}{%
      S_E=\tfrac{P_E}{A_w}=\tfrac{4\pi\,P_E}{\lambda^2\,G}}} &
  \parbox[t]{\linewidth}{%
    \formelblock{Leistungsdichte am Empfänger ($r$)}{%
      S_r=\tfrac{P_S}{4\pi r^2}\quad(\text{Kugelstrahler})}\\[0.2em]%
    \formelblock{Gewinn der Störquelle}{%
      G=\tfrac{S_E}{S_r}}} &
  \parbox[t]{\linewidth}{%
    \infoblock{\textbf{Merke:} Selbst eine \glqq Antenne\grqq,
      die $100\times$ schlechter ist als ein Kugelstrahler
      ($G\!=\!-20$\,dB), kann einen $100$\,m entfernten
      Rundfunkempfänger stören.\\[0.2em]
      Gute Strahler: Drähte/Stäbe, Leitungen gegen Masse,
      Schlitze, Öffnungen $\Rightarrow$ klein/kurz halten!}} \\[0.3em]
\end{formelreihe}
\bereichende


\end{document}
